\documentclass[11pt,a4paper]{book}
\usepackage{fontspec}
\setmainfont{FreeSerif}
\setmonofont{FreeMono}[Scale=0.9]
\newfontfamily\arabicfont[Script=Arabic,Scale=1.12]{DejaVu Sans}
\ifdefined\luatexversion
  % LuaLaTeX: native right-to-left direction (not tested in the build environment)
  \newcommand{\ar}[1]{{\arabicfont\textdir TRT #1\textdir TLT}}
\else
  % XeLaTeX: TeXXeT right-to-left runs
  \TeXXeTstate=1
  \newcommand{\ar}[1]{{\arabicfont\beginR #1\endR}}
\fi
\newcommand{\arq}[1]{\ar{#1}}
\usepackage[margin=2.6cm]{geometry}
\usepackage{booktabs,longtable,enumitem,xcolor}
\usepackage[hidelinks]{hyperref}
\setlist{itemsep=2pt,topsep=4pt}
\setlength{\parindent}{0pt}\setlength{\parskip}{6pt}
\newcommand{\code}[1]{\texttt{#1}}
\newcommand{\NAttested}{45}
\newcommand{\NExact}{33}
\newcommand{\NLetters}{35}
\newcommand{\PctExact}{73}
\newcommand{\PctLetters}{78}
\newcommand{\NDiffers}{10}
\newcommand{\VCensusB}{7}
\newcommand{\VCensusF}{5}
\newcommand{\NFreq}{5000}
\newcommand{\PctSharedFull}{9.8}
\newcommand{\PctSharedSkel}{14.2}
\newcommand{\AtlasExtended}{125}
\newcommand{\RtWords}{2000}
\newcommand{\RtTopOne}{92}
\newcommand{\RtTopFive}{99}
\newcommand{\RtMeanCand}{8.8}
\newcommand{\NApprovedWords}{25}\newcommand{\NApprovedSame}{22}

\title{English in Arabic Script\\[4pt]\large Notation, Evidence and Tools}
\author{Flyxion\\Independent Researcher}
\date{October 2026}

\begin{document}
\frontmatter
\maketitle
\chapter*{Abstract}
This monograph specifies a writing system in which ordinary English is written, word for word and without translation, in standard Arabic letters with optional short vowel marks. It reconstructs the system from the surviving evidence: a handwritten page, the typed passage it copies, and several generated test texts. It records which conventions the evidence settles, which it leaves open, and exposes each open point as a named option in a Python toolkit. The toolkit contains a forward transliterator driven by the CMU Pronouncing Dictionary, a reverse reader that returns ranked candidates, an inventory linter, a corpus statistics tool, and a glossary generator. Every table and figure of agreement in the text is produced by the code and the specimens, so the numbers cannot drift from the implementation.

\tableofcontents
\mainmatter
\chapter{Scope and Definition}

\section{What is being written}
The system writes English. Vocabulary, word order, sentence boundaries and punctuation stay English; only the letters change. A reader who sounds out the Arabic letters recovers the sound of the English word, and a reader who knows English recovers its meaning. Nothing is translated. The sentence
\begin{quote}\ar{إِنْ ذِسْ بَارْتْ أُفْ هِيلِنْ رُولِنْزْ سَايْكُوسِينِمَا}\end{quote}
means \emph{In this part of Helen Rollins' psychocinema} because it sounds like that English phrase, not because any word in it is an Arabic word.

The intended product is a book written in this orthography, accompanied by a pronunciation key and parallel examples. The tools described here exist to produce and check such a book, not to replace the decisions about how the orthography should look.

\section{Constraints}
Four constraints recur in the source material and define the system.
\begin{enumerate}
\item \textbf{English throughout.} Earlier generated texts repeatedly drifted into Arabic vocabulary, for instance spelling out numbers as Arabic number words or substituting a real Arabic word where an English word was expected. Each such drift is an error against this constraint.
\item \textbf{Standard Arabic letters only.} The Persian letters \ar{پ}, \ar{ڤ}, \ar{گ}, \ar{چ}, \ar{ژ} are excluded. English sounds that Arabic lacks are approximated with standard letters.
\item \textbf{Short vowels are optional marks.} Fatha, damma, kasra and sukoon may be written to help the reader. Both fully marked and unmarked presentations were requested, so mark density is a presentation setting, not part of the spelling.
\item \textbf{Approximation is acknowledged.} Several English sounds share a letter, and several English vowels share a representation. The orthography is an approximate phonetic notation. Claims that it reproduces English pronunciation exactly are not made here.
\end{enumerate}

\section{What is not claimed}
Interpretive frameworks attached to the project in earlier discussions, including semiotic or epistemological readings, are outside this monograph. They would need separate justification before becoming part of the foundation. The keyboard layout used for typing the script is treated as an input method, distinct from the sound key (Chapter~\ref{ch:typesetting}).

\section{Organisation}
Part one of the work (Chapters 2 to 5) reconstructs the notation from evidence. Chapters 6 to 10 describe the tools and how they are checked. Chapters 11 to 13 cover teaching formats, typesetting and open work. The appendices give the command line, the full list of attested forms and the legacy sources.

\section{What the notation is for}
The notation serves three uses that pull in slightly different directions. As a reading aid, it lets a reader of Arabic script sound out English text without any knowledge of English spelling, because every word is written as it is said. As a writing system in its own right, it allows English to be handwritten quickly with a small inventory of letters and marks, which is how the handwritten page was produced. As a typographic exercise, it asks how far a script built for another language can carry English before it breaks. The tools in this book support all three. They differ only in how many marks are shown, which is why mark density is a presentation setting and not a change of notation.

\section{Design principles}
Six principles run through the later chapters. They are stated here once so that individual rules can be judged against them.
\begin{enumerate}
\item \textbf{English words, English order.} Nothing is translated and nothing is reordered. A line of the notation has the same number of words as the English line it comes from, and a reader can align the two word by word.
\item \textbf{Standard Arabic letters only.} The letters of the basic Arabic block are used, together with fatha, damma, kasra and sukoon. Persian letters such as \ar{پ}, \ar{چ}, \ar{گ} and \ar{ڤ} are excluded, and near equivalents are accepted instead.
\item \textbf{Near equivalents over new letters.} Where English has a sound that Arabic lacks, the nearest Arabic letter is used. English \emph{p} is written \ar{ب} and English \emph{v} is written \ar{ف}, as a voiced \emph{f}. Exact recovery of short words is not required.
\item \textbf{Marks are optional.} The same letters can be shown with every mark, with vowel marks only, or with none, and the choice never alters a letter.
\item \textbf{Evidence before invention.} A rule is added only when specimens support it, and each unsupported choice is a switch with a stated default.
\item \textbf{Ambiguity is reported.} The notation merges some distinctions, so the reverse tool returns ranked candidates and does not pretend to recover spelling exactly.
\end{enumerate}

\section{How to approach the text}
A reader of this notation is reading English. The Arabic letters are a set of learned conventions for writing it, and they do not turn the text into Arabic or into a transliteration to be decoded each time. At first, sounding out every letter is the natural strategy, and it works because the notation is written as the words are said. With practice the strategy fades. Familiar words become recognisable as whole shapes, as printed English words do, and frequent words such as \ar{ذَا} (\emph{the}) and \ar{أَنْدْ} (\emph{and}) are recognised before their letters are consciously read. Letter-by-letter decoding is therefore a stage in learning, and it is not the reading experience the notation aims at. Later chapters that give a Latin respelling beside the Arabic-script line offer it as a support for that early stage.

\section{A first reading}
The sentences below were composed in the notation and then checked against the toolkit. Each block gives the Arabic-script form with all marks, a plain Latin respelling that is only a reading aid, and the English original. The second line of the first two blocks shows the same words with all marks removed, which is how the notation looks at its lightest.

\begin{quote}\noindent \ar{ذَا بُوكْ إِزْ رِتِنْ إِنْ إِنْغْلِشْ.}\\[2pt]
\ar{ذا بوك إز رتن إن إنغلش.}\\[1pt]
\textit{dhu buk iz ritun in ingglish}\\[1pt]
The book is written in English.
\end{quote}
\begin{quote}\noindent \ar{ذَا لِتَرْزْ تْشَيْنْجْ، بَتْ ذَا وُرْدْزْ رِمَيْنْ.}\\[2pt]
\ar{ذا لترز تشينج، بت ذا وردز رمين.}\\[1pt]
\textit{dhu leterz cheynj but dhu werdz rimeyn}\\[1pt]
The letters change, but the words remain.
\end{quote}
\begin{quote}\noindent \ar{أَ رِيدَرْ لَرْنْزْ ذَا سْكْرِبْتْ أَنْدْ بِغِنْزْ تُو رِيكُغْنَايْزْ أَ فَمِيلْيَرْ فُويْسْ.}\\[2pt]
\ar{أ ريدر لرنز ذا سكربت أند بغنز تو ريكغنايز أ فميلير فويس.}\\[1pt]
\textit{u reeder lernz dhu skript und biginz too rekugnayz u fumilyer voys}\\[1pt]
A reader learns the script and begins to recognise a familiar voice.
\end{quote}

Several choices of the system are visible even in these three sentences. English \emph{p} in \emph{script} is written \ar{ب}. English \emph{v} in \emph{voice} is written \ar{ف}. The \emph{g} of \emph{English} and \emph{begins} is written \ar{غ}, and \emph{ch} in \emph{change} is the two-letter sequence \ar{تْش}. The word \emph{book} keeps a waw because English writes two letters for its vowel, a point that Chapter~\ref{ch:vowels} treats as a general principle.

\section{Plan of the book}
Chapters~2 and~3 describe the specimens and separate what they establish from what remains open. Chapters~4 and~5 state the sound key and the spelling-conditioned vowel rules. Chapters~6 and~7 describe the software architecture and the forward transliterator, and Chapter~8 the reverse reader. Chapters~9 and~10 cover checking and evaluation. Chapter~11 collects worked passages, Chapter~12 deals with typesetting and input, and Chapter~13 lists the open questions. The appendices give the command-line reference, the attested forms and the legacy sources.

\chapter{Specimens and Provenance}

\section{Evidence tiers}
Not all specimens carry equal authority. The reconstruction ranks them.

\begin{description}
\item[Tier 1: the handwritten page.] A page written by hand in the orthography, later transcribed with every mark. It contains four lines:
\begin{quote}\ar{كَابِتَالِسْتْ يُتُوبِيَانِيزْمْ}\\
\ar{ذَا رِبْرِشَنْ أُفْ لَاكْ أَنْدْ إِتْسْ كَلْتْشَرَلْ سِمْبْتُمْزْ.}\\
\ar{إِنْ ذِسْ بَارْتْ أُفْ هِيلِنْ رُولِنْزْ سَايْكُوسِينِمَا،}\\
\ar{شِي دِلْفْزْ إِنْتُو ذَا نَيْتْشَر أُفْ لَاكْ.}\end{quote}
It is the only specimen produced by the author's own hand and is treated as authoritative.
\item[Tier 2: the typed passage.] A longer text of which the page is a copy of the opening. It was produced by an assistant working under the author's instructions and accepted as good, with the explicit acknowledgement of glitches and inconsistencies. Acceptance of a specimen did not certify every spelling in it.
\item[Tier 3: the glossary passage.] An earlier passage with a line-by-line glossary. Its commentary contradicts its examples: it announces one treatment of $p$ and $v$ while the examples show another. Examples outrank commentary.
\item[Tier 4: generated test texts.] Summaries of \emph{Micromegas} and \emph{Atlas Shrugged}, transliterated by an assistant. The \emph{Atlas} text uses four Persian letters; the \emph{Micromegas} text uses one.
\item[Unusable: the chat export.] The original conversation titled \emph{English in Arabic Text} survives only as a text export in which every Arabic character has been replaced by the replacement character U+FFFD. It documents the instructions given, but supplies no letter-level evidence.
\end{description}

\section{Inventory of the specimens}
Table~\ref{tab:inventory} counts letters and marks per specimen using the corpus tool (\code{arabglish stats}) and the linter (\code{arabglish lint}). The page and typed passage use only standard letters. The \emph{Atlas} text accounts for \AtlasExtended{} occurrences of letters outside the standard inventory.

\begin{table}[ht]\centering\small
\begin{tabular}{@{}lrrll@{}}
\toprule
Specimen & Letters & Marks & Extended letters & Lint issues\\
\midrule
Page transcript & 99 & 80 & none & 0\\
Typed passage & 954 & 791 & none & 0\\
Glossary passage & 711 & 545 & none & 3\\
Micromegas & 1417 & 364 & \ar{ڤ\,17} & 21\\
Atlas Shrugged & 1422 & 922 & \ar{ڤ\,57 گ\,30 پ\,34 چ\,4} & 125\\
\bottomrule
\end{tabular}
\caption{Letter, mark and inventory statistics for each specimen. Extended letters are letters outside the standard Arabic block, with their counts.}\label{tab:inventory}
\end{table}

\section{Legacy tools}
Two earlier tools exist and are preserved verbatim in the project under \code{legacy/}.

\textbf{The keyboard layout.} An AutoHotkey script maps Latin keys to Arabic letters one key at a time. Its assignments, such as the key $p$ producing a two-letter sequence, are key positions chosen for typing. They record how to enter characters, not how English sounds are written, and earlier assistant commentary that read them as sound correspondences conflated the two.

\textbf{The Python script.} A pronunciation-driven script uses the CMU Pronouncing Dictionary to obtain English phonemes and maps them to Arabic characters. The preserved copy is the revision with Persian letters removed, as it appears in a pasted conversation, so it may differ from the file that was run. Its fixed word table (\emph{the, this, that, of, and, in, to, a, is}) is adopted in the toolkit. Its structural defects are listed in Appendix~\ref{app:legacy}.

\section{How the specimens were transcribed}
Three of the specimens are the author's own writing. The handwritten page was photographed and typed out letter by letter with every mark, so that the transcript preserves the author's vowel marks and sukoon. The typed passage and the glossary passage were already in digital form. The generated texts, which the author produced with language-model help in earlier sessions, were taken as found.

Transcription from the photograph has two hazards. Fatha, damma and kasra are small and can be confused with dots and with each other, and a sukoon above a letter can look like a damma. Wherever the photograph was ambiguous the typed transcript supplied by the author was treated as authoritative. The transcript is stored in \texttt{specimens/page\_handwritten\_transcript.txt} and its four lines are aligned with the English by \texttt{scripts/build\_attested.py}, which stops with an error if the token counts disagree.

\section{Alignment of Arabic tokens to English words}
Every comparison in this book depends on knowing which Arabic token stands for which English word. For the page and the typed passage the author supplied the English text, and tokens were paired in order. The script asserts that the number of Arabic tokens equals the number of English words in each paragraph, so an insertion or omission anywhere in the text is caught instead of silently shifting every later pair. The result is the list of \NAttested{} word forms in Appendix~\ref{app:attested}.

\section{Strengths and weaknesses of each specimen}
\begin{description}
\item[Handwritten page.] Closest to the author's practice and fully marked, but only four short lines.
\item[Typed passage.] About one thousand letters of fully marked text in the author's conventions. It is the main source of attested forms. It was produced partly by typing and partly with assistance, so it carries some noise, which Chapter~\ref{ch:vowels} notes where it matters.
\item[Glossary passage.] Marks and consonants are consistent, but its own commentary contradicts its letters for the choice of \emph{v}. Its letters were weighted above its commentary.
\item[Micromegas and Atlas.] Early generated texts. They show the system at an earlier stage and with Persian letters in the Atlas text. They are used for the inventory counts and the linter, and not for attested forms.
\end{description}

\section{What the specimens cannot show}
No specimen was written under controlled conditions. The author knew the conventions while writing, and the specimens were not produced from dictated or unseen text. They therefore show how the author writes words the author already has in mind, and give little information about how the author would write a word met for the first time. This limitation returns in Chapter~\ref{ch:eval}, where agreement figures are discussed, and in Chapter~\ref{ch:open}, where held-out specimens are the first request.

\chapter{Evidence and Open Decisions}\label{ch:evidence}

\section{Consonants}
Table~\ref{tab:consonants} gives the consonant key as implemented, with a status for each sound. \emph{Firm} means the specimens agree. \emph{Open} means they conflict. \emph{Proposed} means the author has named a preference that later generated passages did not follow. \emph{Provisional} means no specimen word settles it.

\begin{table}[ht]\centering\small
\begin{tabular}{@{}llcll@{}}
\toprule
ARPAbet & Example & Letters & Status & Note\\
\midrule
P & part & \ar{ب} & firm & ب in every specimen; no specimen uses peh\\
B & big & \ar{ب} & firm & \\
T & top & \ar{ت} & firm & \\
D & dog & \ar{د} & firm & \\
K & cat & \ar{ك} & firm & \\
G & gap & \ar{غ} & proposed & غ proposed by the author; ج appears in generated passages\\
F & fan & \ar{ف} & firm & \\
V & very & \ar{ف} & open & ف by default (v is a voiced f); ب in universal, individual; see census\\
TH & think & \ar{ث} & firm & ث\\
DH & this & \ar{ذ} & firm & ذ\\
S & sun & \ar{س} & firm & \\
Z & zoo & \ar{ز} & firm & \\
SH & she & \ar{ش} & firm & \\
ZH & vision & \ar{ز} & provisional & ز in the cleaned script; no specimen word\\
HH & hat & \ar{ه} & firm & \\
CH & nature & \ar{تش} & firm & تْش on the handwritten page; ش in one typed word\\
JH & subject & \ar{ج} & firm & ج\\
M & man & \ar{م} & firm & \\
N & no & \ar{ن} & firm & \\
NG & sing & \ar{نغ} & provisional & follows G; typed passages show نْجْ\\
L & lack & \ar{ل} & firm & \\
R & red & \ar{ر} & firm & \\
W & way & \ar{و} & firm & و as a consonant\\
Y & yes & \ar{ي} & firm & ي as a consonant\\
\bottomrule
\end{tabular}
\caption{Consonant key (default profile). English $p$ is written \ar{ب} in every specimen.}\label{tab:consonants}
\end{table}

\section{English v}
The evidence on $v$ is divided. Table~\ref{tab:vcensus} lists every $v$ word that could be checked in the specimens. The letter \ar{ب} appears \VCensusB{} times and \ar{ف} appears \VCensusF{} times. The author's own page writes \emph{delves} with \ar{ف}. The typed passage writes \emph{universal}, \emph{individual}, \emph{however}, \emph{generative}, \emph{never}, \emph{resolved} and \emph{deliver} with \ar{ب}. The earlier glossary passage writes \emph{subjective}, \emph{living}, \emph{reveal} and \emph{selves} with \ar{ف}, while its commentary says the opposite.

\begin{table}[ht]\centering\small
\begin{tabular}{@{}llcl@{}}
\toprule
English & Spelling & Letter for v & Specimen\\
\midrule
delves & \ar{دِلْفْزْ} & \ar{ف} & page, typed\\
universal & \ar{يُونِبِرْسَلْ} & \ar{ب} & typed\\
individual & \ar{إِنْدِبِدْيُوَالْ} & \ar{ب} & typed\\
however & \ar{هَوِبَرْ} & \ar{ب} & typed\\
generative & \ar{جِنِرَاتِبْ} & \ar{ب} & typed\\
never & \ar{نِبَرْ} & \ar{ب} & typed\\
resolved & \ar{رِزُولْبْدْ} & \ar{ب} & typed\\
deliver & \ar{دِلِبَرْ} & \ar{ب} & typed\\
subjective & \ar{سُوبْجِكْتِفْ} & \ar{ف} & glossary\\
living & \ar{لَافِنْجْ} & \ar{ف} & glossary\\
reveal & \ar{رِيفِيلْ} & \ar{ف} & glossary\\
selves & \ar{إِنْسْلِيفْزْ} & \ar{ف} & glossary\\
\bottomrule
\end{tabular}
\caption{Census of English $v$ in the specimens that could be aligned.}\label{tab:vcensus}
\end{table}

The earlier Python script used a voicing heuristic, choosing \ar{ب} between voiced neighbours and \ar{ف} otherwise. That heuristic cannot reproduce the census: \emph{delves} and \emph{resolved} have identical neighbourhoods (a liquid before $v$, a voiced consonant after) yet receive different letters. The toolkit therefore offers three policies, \code{b}, \code{f} and \code{context}, and a per-word override through the attested list. The default is \ar{ف}, on the author's reasoning that $v$ is a voiced $f$; this also matches the old program and the page's \emph{delves}. The typed passage's \emph{individual} and \emph{universal} use \ar{ب}, and \code{--v b} reproduces that choice.

\section{Hard g and the velar nasal}
The author proposed \ar{غ} (ghain) for hard $g$. Generated passages used \ar{ج} (jeem), and the cleaned Python script mapped $g$ to \ar{ج}, which merges $g$ with $j$ in \emph{subject}. The default is \ar{غ}, with \code{--g jeem} available. The velar nasal follows $g$: \ar{نغ} by default and \ar{نج} under \code{jeem}.

\section{The affricate ch}
The page writes \emph{cultural} and \emph{nature} with the two-letter sequence \ar{تْش}, taking sukoon on the first letter. One typed word, \emph{nature}, uses \ar{ش} alone. The sequence is the default. When the sequence closes a word in full marks both letters carry sukoon, as in \ar{وِتْشْ} \emph{which}.

\section{Remaining decisions}
Table~\ref{tab:decisions} lists the decisions and where the toolkit exposes them.

\begin{table}[ht]\centering\small
\begin{tabular}{@{}p{3.3cm}p{4.6cm}p{5.2cm}@{}}
\toprule
Decision & Default & Switch\\
\midrule
English $v$ & \ar{ف} & \code{--v b|f|context}\\
Hard $g$ & \ar{غ} & \code{--g ghain|jeem}\\
Persian letters & excluded & \code{--persian} (for conversion only)\\
Mark density & full marks & \code{--density full|medial|vowels|none}\\
\emph{the} before a vowel & \ar{ذَا} always & \code{--the-before-vowel}\\
Numbers & Arabic-Indic digits & \code{--digits indic|latin|spoken}\\
Specimen forms & computed forms & \code{--attested}\\
Doubled letters & not marked (no shadda) & none\\
\bottomrule
\end{tabular}
\caption{Open decisions and their switches.}\label{tab:decisions}
\end{table}

\section{Decision log}
The consonant and vowel choices are recorded here with the evidence and reasoning behind each, in the order in which they were settled.

\begin{description}
\item[English p.] Written \ar{ب} in every specimen and in the old program. The author states that \ar{ب} for \emph{p} is acceptable and that exact recovery of short words is not required. Settled.
\item[English v.] The old program wrote \ar{ف}. The typed passage writes \ar{ب} in \emph{universal} and \emph{individual} and the page writes \ar{ف} in \emph{delves}. The author's reasoning is that \emph{v} is only a voiced \emph{f}, so the default is \ar{ف}. The alternative is kept as \code{--v b}.
\item[Hard g.] The old program wrote \ar{ج}, merging it with \emph{j}. The author later proposed \ar{غ}, which keeps \emph{g} distinct from \emph{j}. The default is \ar{غ} and \code{--g jeem} restores the older choice.
\item[English ch.] The page writes the two-letter sequence \ar{تْش}, the typed passage writes \ar{ش} in \emph{nature}, and the old program wrote \ar{ش}. The sequence is the default because it keeps \emph{ch} separate from \emph{sh}.
\item[English ng.] Written as nasal plus ghain, \ar{نغ}. The old program wrote \ar{ن} alone. A following hard \emph{g}, as in \emph{English} and \emph{finger}, is not written twice.
\item[Voiced th and voiceless th.] Written \ar{ذ} and \ar{ث} in every source.
\item[Zh.] Written \ar{ز}. There is no better standard Arabic letter and no specimen word with the sound is attested.
\item[Vowel colour of unstressed syllables.] Taken from the vowel letter of the spelling, with the exceptions listed in Chapter~\ref{ch:vowels}.
\end{description}

\section{Consonant examples}
The following words exercise the consonant decisions: \emph{thing}, \emph{think}, \emph{this}, \emph{that}, \emph{church}, \emph{judge}, \emph{measure}, \emph{vision}, \emph{singer} and \emph{finger}, with a few clusters.

\begin{table}[ht]\centering\small
\begin{tabular}{@{}lll@{}}
\toprule
English & Arabic script & Respelling\\
\midrule
thing & \ar{ثِنْغْ} & \textit{thing}\\
think & \ar{ثِنْغْكْ} & \textit{thingk}\\
this & \ar{ذِسْ} & \textit{dhis}\\
that & \ar{ذَتْ} & \textit{dhat}\\
church & \ar{تْشَرْتْشْ} & \textit{cherch}\\
change & \ar{تْشَيْنْجْ} & \textit{cheynj}\\
judge & \ar{جَجْ} & \textit{juj}\\
measure & \ar{مِيزَرْ} & \textit{mezher}\\
vision & \ar{فِيزَنْ} & \textit{vizhun}\\
singer & \ar{سِنْغَرْ} & \textit{singer}\\
finger & \ar{فِنْغَرْ} & \textit{fingger}\\
garden & \ar{غَارْدِنْ} & \textit{gaardun}\\
\bottomrule
\end{tabular}
\caption{Words that exercise the consonant decisions.}\label{tab:wcons}
\end{table}

\section{Decisions that need no more evidence}
Not every open choice is equally pressing. The choice for \emph{v} affects a small group of words and is a single switch. The choice for hard \emph{g} affects more words but the two options are equally readable. What matters most for the next stage is a larger set of fully marked words in the author's own hand, because most of the vowel rules rest on a few dozen forms.

\chapter{The Sound Key}

\section{Vowels}
Each English vowel is written as a short mark on the preceding consonant and, for long vowels and glides, a vowel letter. Table~\ref{tab:vowels} shows the default correspondences using words from the page. Phonemes are named in ARPAbet, the notation of the pronouncing dictionary.

\begin{table}[ht]\centering\small
\begin{tabular}{@{}lllll@{}}
\toprule
ARPAbet & Mark & Letters & Example & Rendered\\
\midrule
AA & fatha & \ar{ا} & part & \ar{بَارْتْ}\\
AE & fatha & -- & lack & \ar{لَاكْ}\\
AH & fatha & -- & cultural & \ar{كَلْتْشَرَلْ}\\
AO & damma & \ar{و} & on & \ar{أُونْ}\\
AW & fatha & \ar{او} & flower & \ar{فْلَاوَرْ}\\
AY & fatha & \ar{اي} & psychocinema & \ar{سَايْكُوسِينِمَا}\\
EH & kasra & -- & delves & \ar{دِلْفْزْ}\\
ER & fatha & \ar{ر} & nature & \ar{نَيْتْشَرْ}\\
EY & fatha & \ar{ي} & nature & \ar{نَيْتْشَرْ}\\
IH & kasra & -- & this & \ar{ذِسْ}\\
IY & kasra & \ar{ي} & she & \ar{شِي}\\
OW & damma & \ar{و} & Rollins & \ar{رُولِنْزْ}\\
OY & damma & \ar{وي} & void & \ar{فُويْدْ}\\
UH & damma & -- & could & \ar{كُودْ}\\
UW & damma & \ar{و} & into & \ar{إِنْتُو}\\
\bottomrule
\end{tabular}
\caption{Vowel key. The mark is the short vowel; the letters are the vowel letters added for long vowels and glides.}\label{tab:vowels}
\end{table}

Three kinds of vowel letter occur. A \emph{long} letter, as in \ar{سِينِمَا}, takes no sukoon. A \emph{glide} letter, the second element of a diphthong as in \ar{سَايْكُو}, takes sukoon. A \emph{rhotic} vowel is written as a mark followed by \ar{ر}, as in \ar{شَرْ}.

\section{Sukoon}
Sukoon marks a consonant followed by no vowel. In full density every such consonant carries it, including the last consonant of a word, as in \ar{لَاكْ} \emph{lack} and \ar{كَابِتَالِسْتْ} \emph{capitalist}. Sukoon records the absence of a vowel; it does not make a letter silent, and English double letters are not marked with shadda.

\section{Word-initial vowels}
A word that begins with a vowel takes a hamza seat: \ar{أ} for fatha and damma, \ar{إ} for kasra, as in \ar{أُفْ} \emph{of} and \ar{إِنْ} \emph{in}. An initial diphthong beginning with fatha and alef is written with \ar{آ}. A vowel that follows another vowel is hosted by a preceding \ar{ي} when there is one (\ar{ثِيُ}) and otherwise takes a seat.

\section{Frequent words}
Frequent function words have fixed spellings (Table~\ref{tab:common}). These come from the original fixed word table and from forms that the page and the typed passage share. They also resolve the contextual variation of \emph{the} and \emph{a}: the toolkit writes \ar{ذَا} before a vowel too, as the typed passage does, and offers the phonemic alternative \ar{ذِي} as an option.

\begin{table}[ht]\centering\small
\begin{tabular}{@{}ll@{\qquad}ll@{}}
\toprule
English & Spelling & English & Spelling\\
\midrule
the & \ar{ذَا} & is & \ar{إِزْ}\\
this & \ar{ذِسْ} & she & \ar{شِي}\\
that & \ar{ذَتْ} & it & \ar{إِتْ}\\
of & \ar{أُفْ} & its & \ar{إِتْسْ}\\
and & \ar{أَنْدْ} & as & \ar{أَزْ}\\
in & \ar{إِنْ} & into & \ar{إِنْتُو}\\
to & \ar{تُو} & with & \ar{وِذْ}\\
a & \ar{أَ} & are & \ar{أَرْ}\\
\bottomrule
\end{tabular}
\caption{Fixed spellings of frequent words.}\label{tab:common}
\end{table}

\section{Punctuation and numbers}
The comma, semicolon and question mark take their Arabic forms \ar{،}, \ar{؛}, \ar{؟}. A possessive apostrophe after $s$ is silent, as on the page. Digits default to Arabic-Indic numerals. A \code{spoken} option writes numbers as English number words in the notation (\emph{twenty five}), which avoids the earlier failure of rendering numbers as Arabic number words.

\section{Letter sound and assigned function}
Each Arabic letter has an ordinary sound in Arabic, and in this notation it also has an assigned function in writing English. The two often coincide, as with \ar{ب} for \emph{b} and \ar{ل} for \emph{l}, and sometimes they differ. The letter \ar{غ} is a voiced uvular fricative in Arabic, and it is used here for the hard English \emph{g} by convention. The same holds for \ar{ث} and \ar{ذ}, which are close to English \emph{th} and \emph{dh} sounds, and for \ar{ب}, which stands for both \emph{b} and \emph{p}. The letters \ar{و} and \ar{ي} serve as consonants (\emph{w} and \emph{y}) and also as vowel letters.

Knowing Arabic script is therefore a useful starting point, since the letter shapes, the direction of writing and the role of the marks are already familiar. It does not supply every reading rule. A reader who pronounces \ar{غ} as in Arabic will mispronounce \emph{good} and \emph{begin}, and a reader who expects \ar{ب} to be only \emph{b} will miss the \emph{p} in \emph{script}. The tables in this chapter give the assigned function, and the Arabic sound is only a mnemonic where it helps.

\section{Short vowels in practice}
Short vowels are written with a single mark on the preceding consonant and no vowel letter. The words below show the five English short vowels, each followed by one consonant or a cluster.

\begin{table}[ht]\centering\small
\begin{tabular}{@{}lll@{}}
\toprule
English & Arabic script & Respelling\\
\midrule
cat & \ar{كَاتْ} & \textit{kat}\\
bed & \ar{بِدْ} & \textit{bed}\\
sit & \ar{سِتْ} & \textit{sit}\\
hot & \ar{هُوتْ} & \textit{haat}\\
cut & \ar{كَتْ} & \textit{kut}\\
put & \ar{بُتْ} & \textit{put}\\
pen & \ar{بِنْ} & \textit{pen}\\
map & \ar{مَابْ} & \textit{map}\\
top & \ar{تُوبْ} & \textit{taap}\\
sun & \ar{سَنْ} & \textit{sun}\\
\bottomrule
\end{tabular}
\caption{Short vowels, written with a mark only.}\label{tab:wshort}
\end{table}

\section{Long vowels and the final e}
English marks long vowels in several ways: a doubled letter, a vowel digraph, or a final silent \emph{e}. The dictionary reading already gives the long vowel in each case, so the notation writes a vowel letter after the mark. The final \emph{e} itself is silent and is not written.

\begin{table}[ht]\centering\small
\begin{tabular}{@{}lll@{}}
\toprule
English & Arabic script & Respelling\\
\midrule
make & \ar{مَيْكْ} & \textit{meyk}\\
time & \ar{تَايْمْ} & \textit{taym}\\
home & \ar{هُومْ} & \textit{hohm}\\
use & \ar{يُوسْ} & \textit{yoos}\\
cake & \ar{كَيْكْ} & \textit{keyk}\\
like & \ar{لَايْكْ} & \textit{layk}\\
phone & \ar{فُونْ} & \textit{fohn}\\
tune & \ar{تُونْ} & \textit{toon}\\
\bottomrule
\end{tabular}
\caption{Long vowels and glides written for the silent-e pattern.}\label{tab:wlong}
\end{table}

\section{Doubled letters and vowel digraphs}
When the English spelling uses two vowel letters for one vowel, the notation keeps a waw or a yeh. This applies to the short vowel of \emph{book} as much as to the long vowel of \emph{food}. The spelling carries the weight here and the Arabic letter signals it. Chapter~\ref{ch:vowels} explains why \emph{book} is \ar{بُوكْ} and not \ar{بُكْ}.

\begin{table}[ht]\centering\small
\begin{tabular}{@{}lll@{}}
\toprule
English & Arabic script & Respelling\\
\midrule
book & \ar{بُوكْ} & \textit{buk}\\
look & \ar{لُوكْ} & \textit{luk}\\
good & \ar{غُودْ} & \textit{gud}\\
foot & \ar{فُوتْ} & \textit{fut}\\
wood & \ar{وُودْ} & \textit{wud}\\
food & \ar{فُودْ} & \textit{food}\\
moon & \ar{مُونْ} & \textit{moon}\\
school & \ar{سْكُولْ} & \textit{skool}\\
too & \ar{تُو} & \textit{too}\\
two & \ar{تُو} & \textit{too}\\
blue & \ar{بْلُو} & \textit{bloo}\\
true & \ar{تْرُو} & \textit{troo}\\
\bottomrule
\end{tabular}
\caption{Words with a written double o, and related spellings.}\label{tab:woo}
\end{table}

\begin{table}[ht]\centering\small
\begin{tabular}{@{}lll@{}}
\toprule
English & Arabic script & Respelling\\
\midrule
see & \ar{سِي} & \textit{see}\\
tree & \ar{تْرِي} & \textit{tree}\\
green & \ar{غْرِينْ} & \textit{green}\\
sleep & \ar{سْلِيبْ} & \textit{sleep}\\
head & \ar{هِدْ} & \textit{hed}\\
bread & \ar{بْرِدْ} & \textit{bred}\\
read & \ar{رِدْ} & \textit{red}\\
meet & \ar{مِيتْ} & \textit{meet}\\
meat & \ar{مِيتْ} & \textit{meet}\\
dream & \ar{دْرِيمْ} & \textit{dreem}\\
dead & \ar{دِدْ} & \textit{ded}\\
\bottomrule
\end{tabular}
\caption{Words with ee and ea.}\label{tab:wee}
\end{table}

\section{Glides and diphthongs}
The diphthongs of English are written as a vowel mark followed by a glide letter that takes sukoon, so \emph{day} is \ar{دَيْ}, \emph{light} is \ar{لَايْتْ} and \emph{voice} is \ar{فُويْسْ}. This is the placement seen on the handwritten page.

\begin{table}[ht]\centering\small
\begin{tabular}{@{}lll@{}}
\toprule
English & Arabic script & Respelling\\
\midrule
rain & \ar{رَيْنْ} & \textit{reyn}\\
day & \ar{دَيْ} & \textit{dey}\\
light & \ar{لَايْتْ} & \textit{layt}\\
night & \ar{نَايْتْ} & \textit{nayt}\\
voice & \ar{فُويْسْ} & \textit{voys}\\
flower & \ar{فْلَاوَرْ} & \textit{flower}\\
hour & \ar{آوَرْ} & \textit{ower}\\
boy & \ar{بُويْ} & \textit{boy}\\
toy & \ar{تُويْ} & \textit{toy}\\
noise & \ar{نُويْزْ} & \textit{noyz}\\
now & \ar{نَاوْ} & \textit{now}\\
how & \ar{هَاوْ} & \textit{how}\\
\bottomrule
\end{tabular}
\caption{Glides and diphthongs.}\label{tab:wai}
\end{table}

\section{R-coloured vowels}
A vowel followed by \emph{r} in the dictionary is written as a vowel mark and \ar{ر}. The mark of an unstressed final syllable is fatha in every attested form, as in \ar{نَيْتْشَر} and \ar{هَوِبَرْ}. Stressed syllables take their mark from the spelling.

\begin{table}[ht]\centering\small
\begin{tabular}{@{}lll@{}}
\toprule
English & Arabic script & Respelling\\
\midrule
car & \ar{كَارْ} & \textit{kaar}\\
far & \ar{فَارْ} & \textit{faar}\\
her & \ar{هِرْ} & \textit{her}\\
bird & \ar{بِرْدْ} & \textit{berd}\\
turn & \ar{تَرْنْ} & \textit{tern}\\
word & \ar{وُرْدْ} & \textit{werd}\\
door & \ar{دُورْ} & \textit{dawr}\\
more & \ar{مُورْ} & \textit{mawr}\\
fire & \ar{فَايَرْ} & \textit{fayer}\\
tire & \ar{تَايَرْ} & \textit{tayer}\\
\bottomrule
\end{tabular}
\caption{R-coloured vowels.}\label{tab:wr}
\end{table}

\section{Clusters and sukoon}
English has many consonant clusters that Arabic does not allow at the start of a word. The notation writes them with sukoon on all but the last consonant, so the cluster \emph{str} in \emph{street} appears as \ar{سْتْرِيتْ}. Word-final consonants also take sukoon in the full density. The \texttt{medial} density removes the final sukoon, which leaves clusters inside words intact.

\begin{table}[ht]\centering\small
\begin{tabular}{@{}lll@{}}
\toprule
English & Arabic script & Respelling\\
\midrule
string & \ar{سْتْرِنْغْ} & \textit{string}\\
spring & \ar{سْبْرِنْغْ} & \textit{spring}\\
splash & \ar{سْبْلَاشْ} & \textit{splash}\\
street & \ar{سْتْرِيتْ} & \textit{street}\\
strength & \ar{سْتْرِنْغْكْثْ} & \textit{strengkth}\\
twelfth & \ar{تْوِلْفْثْ} & \textit{twelfth}\\
asked & \ar{أَسْكْتْ} & \textit{askt}\\
texts & \ar{تِكْسْتْسْ} & \textit{teksts}\\
\bottomrule
\end{tabular}
\caption{Clusters, written with sukoon.}\label{tab:wclusters}
\end{table}

\chapter{Vowels Conditioned by Spelling}\label{ch:vowels}

\section{Observation}
A purely phonemic mapping from the pronouncing dictionary reproduces the page's consonants well and its vowels less well. Comparing the computed forms with the attested forms shows a regular departure: the letter used for a vowel sometimes follows the English spelling rather than the sound.

\begin{enumerate}
\item \textbf{Stressed o.} The dictionary gives \emph{Rollins}, \emph{on}, \emph{psychology} and \emph{concept} an open back vowel. The specimens write them with damma and waw: \ar{رُولِنْزْ}, \ar{أُونْ}, \ar{سَايْكُولُوجِي}, \ar{كُونْسِبْتْ}. Contrast \emph{part}, written \ar{بَارْتْ} with fatha and alef, where the spelling letter is $a$.
\item \textbf{Open-syllable e and i.} \emph{Helen} and \emph{psychocinema} are written with kasra and \ar{ي}: \ar{هِيلِنْ}, \ar{سِينِمَا}. The stressed vowel letter is followed by one consonant and another vowel. In \emph{delves}, \emph{this} and \emph{its}, where the vowel letter is followed by a cluster, only kasra is written.
\item \textbf{Schwa colour.} An unstressed vowel takes the colour of its letter: $a$ gives fatha, $e$ and $i$ give kasra, $o$ gives damma, as in \ar{سِمْبْتُمْزْ} \emph{symptoms}.
\item \textbf{Short initial vowels.} An unstressed vowel in the first syllable is written short, with no vowel letter: \ar{رِبْرِشَنْ} \emph{repression}, \ar{يُتُوبِيَانِيزْمْ} \emph{utopianism}.
\item \textbf{Hosted hiatus.} A vowel after \ar{ي} is written on the \ar{ي}, as in \ar{بِيَ} and \ar{ثِيُ}.
\end{enumerate}

\section{Why pronunciation and spelling both contribute}
The notation could have been defined as a pure phonetic transcription, with each vowel written from its dictionary sound alone. The author's practice shows that it is not. Familiar English words carry a visual identity that readers use in recognising them, and the notation keeps some of it. The word \emph{book} is written \ar{بُوكْ} with a waw because English writes two vowel letters, and the waw retains that feature even though it does not reproduce the short vowel narrowly. The waw is a deliberate orthographic choice. It is not an error of pronunciation to be corrected.

The consequence for later revision is practical. A departure from dictionary pronunciation in an attested form is not automatically a mistake. It may be the spelling showing through, and a rule that removes it would make the notation more phonetic and less like the author's writing. Before such a form is changed, the question to ask is whether the departure is systematic across words with the same spelling, as the waw of \emph{book}, \emph{good} and \emph{foot} is, or whether it is a lapse.

\section{Reading of the pattern}
These observations suggest that the notation is a hybrid. It is phonetic in its consonants and in the broad shape of its syllables, and orthographic in the choice of some vowel letters. The evidence is one author's page, a typed passage and generated texts, so the hybrid reading is a hypothesis about how the notation behaves, not a finding about how it must behave. The rules implemented from it are listed with their results in Chapter~\ref{ch:eval}.

\section{Remaining disagreements}
After the rules are applied, \NDiffers{} of the \NAttested{} attested forms still differ from the computed forms by more than marks. Table~\ref{tab:differs} lists them. They fall into five groups:
\begin{description}
\item[Lexical alef.] \emph{capitalist}, \emph{capitalism}, \emph{utopianism} and \emph{psychoanalytic} carry an alef where the rules give only fatha. Other words with the same spelling pattern, such as \emph{fundamental} and \emph{cultural}, do not. No rule from spelling or stress separates them, so these stay in the attested list.
\item[The choice of v.] \emph{individual} and \emph{universal}, written with \ar{ب} in the typed passage, discussed in Chapter~\ref{ch:evidence}. The default \ar{ف} is the author's decision and the specimens disagree with it in these two words.
\item[Failed alignment.] When the number of vowel letter groups differs from the number of vowel phonemes, as in \emph{ideology} and \emph{individual}, spelling-based colouring is skipped and the phonemic default applies.
\item[Seat letter.] The article \emph{a} is written \ar{اَ} in the typed passage and \ar{أَ} in the original fixed table. Both appear in the sources; the toolkit keeps the hamza form.
\item[Generated-text noise.] \emph{nature} in the typed passage, which uses \ar{ش} where the page has \ar{تْش}, and \emph{structures}, whose typed spelling departs from both the page conventions and the sound.
\end{description}

\begin{table}[ht]\centering\small
\begin{tabular}{@{}lll@{}}
\toprule
English & Attested & Computed\\
\midrule
a & \ar{اَ} & \ar{أَ}\\
capitalism & \ar{كَابِتَالِيزْمْ} & \ar{كَابِتَلِزْمْ}\\
capitalist & \ar{كَابِتَالِسْتْ} & \ar{كَابِتَلِسْتْ}\\
ideology & \ar{آيدِيُولُوجِي} & \ar{آيْدِيَالَجِي}\\
individual & \ar{إِنْدِبِدْيُوَالْ} & \ar{إِنْدَفِجَوَلْ}\\
nature & \ar{نَيْشَرْ} & \ar{نَيْتْشَرْ}\\
psychoanalytic & \ar{سَايْكُوأَنَالِيتِيكْ} & \ar{سَايْكُوأَنَلِتِكْ}\\
structures & \ar{سْتْرَكْشُورْزْ} & \ar{سْتْرَكْتْشَرْزْ}\\
universal & \ar{يُونِبِرْسَلْ} & \ar{يُونِفِرْسَلْ}\\
utopianism & \ar{يُتُوبِيَانِيزْمْ} & \ar{يُتُوبِيَنِزْمْ}\\
\bottomrule
\end{tabular}
\caption{Attested forms that differ from the computed forms in their letters.}\label{tab:differs}
\end{table}

\section{The rules in order}
The vowel rules are applied in a fixed order for each vowel phoneme, and the first rule that fires decides the mark and the vowel letters.
\begin{enumerate}
\item Schwa (AH) takes the colour of its spelling letter: $e$, $i$ and $y$ give kasra, $o$ gives damma, $a$ gives fatha, with fatha also for the vowel pairs \emph{ia} and \emph{io} and for word-final \emph{-em}. A stressed AH spelled with $o$ is written with damma and waw.
\item A rhotic vowel (ER) is written with fatha when unstressed, with the exception of the prefix \emph{inter-}, which keeps kasra. Stressed ER takes its mark from the spelling letter, but \emph{ea} gives fatha (\emph{learn}).
\item The vowels AA and AO, when the spelling letter is $o$, are written with damma and waw (\emph{on}, \emph{Rollins}).
\item A stressed AE that is not word-initial gains an alef (\emph{lack}).
\item A stressed EH or IH in an open syllable spelled with $e$ or $i$ gains a yeh (\emph{Helen}, \emph{cinema}).
\item An unstressed IY or UW in the first syllable is written short.
\item The vowel UH spelled with a written \emph{oo} or \emph{ou} keeps a waw (\emph{book}, \emph{good}, \emph{foot}).
\end{enumerate}
Each rule can be traced to specimen words, and each has exceptions among the attested forms that are listed in Table~\ref{tab:differs}.

\section{Double o and the word book}
The last rule came from a correction by the author. A composed example wrote \emph{book} as \ar{بُكْ}, following the sound, and the author preferred \ar{بُوكْ} on the ground that English writes two vowel letters and the notation should show them. The author also noted that without the waw an Arabic reader might hear a different vowel. The same reasoning applies to \emph{look}, \emph{good}, \emph{foot} and \emph{wood}, and the toolkit now writes all of them with a waw. It does not yet extend to every doubled vowel, because the specimens contain no attested case of other doubled letters, apart from \emph{ee}, which already receives a yeh as a long vowel.

\section{Hosted vowels after a glide}
When a vowel follows a glide, the mark is placed on the glide, so no separate hamza seat is needed. This accounts for \emph{flower}, \emph{power} and \emph{hour}, where the diphthong in the first syllable ends in a waw glide and the unstressed syllable takes its fatha or kasra on that waw. The same device was already used after \ar{ي}, as in \emph{utopianism}.

\section{Ng and the hard g}
In the dictionary, \emph{finger}, \emph{anger} and \emph{English} have an ng sound followed by a hard g. Written as separate sounds this gives a doubled ghain, which no specimen shows. The toolkit therefore writes the nasal and one ghain, as in \ar{إِنْغْلِشْ}. Words such as \emph{singer} and \emph{longer}, where the dictionary has only the nasal, are written the same way, which is a merger the notation accepts.

\section{Examples}
Table~\ref{tab:whybrid} gives words from the specimens and words that were not in them, written by the rules. The specimen words reproduce their attested forms. The others show how the rules extend.

\begin{table}[ht]\centering\small
\begin{tabular}{@{}lll@{}}
\toprule
English & Arabic script & Respelling\\
\midrule
psychology & \ar{سَايْكُولُوجِي} & \textit{saykaalujee}\\
concept & \ar{كُونْسِبْتْ} & \textit{kaansept}\\
Rollins & \ar{رُولِنْزْ} & \textit{raalinz}\\
on & \ar{أُونْ} & \textit{aan}\\
cinema & \ar{سِينِمَا} & \textit{sinumu}\\
Helen & \ar{هِيلِنْ} & \textit{helun}\\
repression & \ar{رِبْرِشَنْ} & \textit{reepreshun}\\
social & \ar{سُوشَلْ} & \textit{sohshul}\\
nation & \ar{نَيْشَنْ} & \textit{neyshun}\\
system & \ar{سِسْتَمْ} & \textit{sistum}\\
lack & \ar{لَاكْ} & \textit{lak}\\
part & \ar{بَارْتْ} & \textit{paart}\\
flower & \ar{فْلَاوَرْ} & \textit{flower}\\
hour & \ar{آوَرْ} & \textit{ower}\\
power & \ar{بَاوَرْ} & \textit{power}\\
tower & \ar{تَاوَرْ} & \textit{tower}\\
finger & \ar{فِنْغَرْ} & \textit{fingger}\\
English & \ar{إِنْغْلِشْ} & \textit{ingglish}\\
\bottomrule
\end{tabular}
\caption{Words that exercise the spelling-conditioned rules.}\label{tab:whybrid}
\end{table}

\chapter{Architecture of the Toolkit}

\section{Modules}
The toolkit is a Python package, \code{arabglish}, with one module per concern.

\begin{description}
\item[\code{tables}] Sound-to-letter tables and the profile object. Every open decision is a parameter of \code{make\_profile}.
\item[\code{lexicon}] Pronunciation lookup and alignment of vowel phonemes to vowel letters.
\item[\code{g2p}] A small spelling-to-sound fallback for words the dictionary lacks.
\item[\code{translit}] The forward transliterator, mark density, numbers and attested forms.
\item[\code{reader}] The reverse reader and round-trip measurement.
\item[\code{lint}, \code{corpus}] Inventory checking and specimen statistics.
\item[\code{guide}] The Latin respelling used in teaching glossaries.
\item[\code{cli}] The command line.
\end{description}

The command \code{arabglish} offers the subcommands \code{fwd}, \code{rev}, \code{lint}, \code{stats}, \code{glossary} and \code{roundtrip} (Appendix~\ref{app:cli}). The original wrapper script, \code{main.py}, runs unchanged through a small compatibility module.

\section{Pronunciation source}
Pronunciations come from the CMU Pronouncing Dictionary through the \code{cmudict} package, which holds 126,052 entries of General American English. Lookup proceeds in order:
\begin{enumerate}
\item an override file, \code{pron\_overrides.tsv}, for words the dictionary lacks or where a specimen needs a specific reading (\emph{psychocinema}, \emph{psychoanalytic}, \emph{utopianism}, \emph{disemvoweling});
\item the dictionary, using its first pronunciation, and for the reverse index all pronunciations;
\item possessives built from the stem with the regular suffix;
\item the spelling fallback, which logs a warning so that weak spots are visible.
\end{enumerate}

\section{Replacing the pronunciation source}
The notation is tied to one accent by this choice: General American, with its vowels and rhotics. A different source changes the Arabic output for many words. Pronunciation lookup is isolated in the single function \code{pronounce}, which returns a phoneme list and a source label, so another dictionary can be substituted without touching the rest. Candidates include the eSpeak NG phonemizer, which also covers unknown words and offers British voices, and the IPA transcriptions in Wiktionary. Which accent the book should follow is an open question for the author (Chapter~\ref{ch:open}).

\section{Determinism}
For a given option set the output is deterministic. Mark density is applied as a final filter over the fully marked form, so changing density never changes a letter: this is checked by a test over a sample text.

\section{Module by module}
\begin{description}
\item[\texttt{tables.py}] The sound key as data: consonant letters, vowel marks and letters, fixed spellings of frequent words, punctuation, Arabic-Indic digits and the \texttt{Profile} object that records the open decisions.
\item[\texttt{g2p.py}] A small spelling-to-phoneme fallback. It handles the common English digraphs, the magic-e pattern and a few suffixes, and it assigns primary stress to the first vowel. It is deliberately simple, because every word that reaches it is logged and can be added to the override file.
\item[\texttt{lexicon.py}] The pronunciation layer. For a word it tries the override file, then the CMU dictionary, then a possessive rule, then the fallback, and it returns the phonemes with the name of the source. It also aligns the vowel letter groups of the spelling to the vowel phonemes.
\item[\texttt{translit.py}] The forward transliterator: segments, vowel placement, sukoon, density, attested forms, digits and text tokens.
\item[\texttt{reader.py}] The reverse reader and its index.
\item[\texttt{lint.py}] The inventory checker and the extended-letter repair.
\item[\texttt{corpus.py}] Statistics over specimens: letter counts, mark counts and inventory.
\item[\texttt{guide.py}] The Latin respelling used for teaching rows.
\item[\texttt{cli.py}] The command line.
\end{description}

\section{Data files}
Two small tab-separated files sit beside the code. \texttt{pron\_overrides.tsv} lists words whose dictionary pronunciation is replaced, mainly the author's coinages and technical words (\emph{psychocinema}, \emph{disemvoweling}). \texttt{attested.tsv} records word forms taken from the specimens, with the source of each. Both are plain text, so they can be edited by hand and reviewed in a diff.

\section{Pronunciation source}
All pronunciation lookups go through a single function, \texttt{pronounce}. The dictionary is read once and cached. Replacing it with another source, such as a British dictionary, a machine-readable dictionary in a different phone set, or a table of the author's own readings, means changing that function and a mapping from the new phone set to the ARPABET symbols used by the tables. Nothing downstream needs to change.

\section{Diagnostics}
The fallback logs every word it handles, and the unknown-phoneme path logs a warning, so weak spots are visible. The attested-form comparison is a script, not a test, so that it can be rerun and read. The test suite then fixes the numbers it produces, so regressions fail loudly.

\section{Testing}
The test suite has three kinds of test. Unit tests check single rules on single words, such as the waw of \emph{book} and the silent possessive apostrophe in \emph{Rollins'}. Invariant tests check properties that must always hold, such as that changing density never changes the unmarked skeleton. A ratchet test checks the agreement with the attested forms against fixed lower bounds, which are raised whenever the rules improve.

\chapter{The Forward Transliterator}

\section{From words to segments}
For each word the engine obtains a phoneme list, aligns vowel phonemes to the vowel letter groups of the spelling when their counts agree, and builds a list of segments. Each segment is a letter with a kind (consonant, long vowel letter, glide, seat) and an optional mark.

Consonant phonemes append their letters, digraphs such as \ar{تْش} appending two segments of one unit. Vowel phonemes attach their mark to the last consonant, or to a preceding \ar{ي}, or to a new hamza seat, and then append their vowel letters. Rhotic vowels append \ar{ر}.

\section{The sukoon pass}
After all segments are placed, every consonant or glide without a vowel receives sukoon. Long vowel letters and seats do not. The resulting form is the full-density spelling.

\section{Mark density}
A final filter produces the other densities, so density can never change a letter.
\begin{description}
\item[full] all marks, including sukoon on every vowelless consonant.
\item[medial] sukoon removed from the end of the word.
\item[vowels] fatha, damma and kasra only.
\item[none] no marks. This is the consonantal skeleton with vowel letters, the form requested in the earliest experiments.
\end{description}

\section{Attested forms}
With \code{--attested} the engine prefers a form recorded in a specimen whenever the word occurs there. The list is built by aligning each specimen token to its English word and asserting that token counts match (\code{scripts/build\_attested.py}). Page forms are listed before typed forms.

\section{Demonstration}
Table~\ref{tab:demo} shows one sentence of the page under several settings. The first row is the page itself.

\begin{table}[ht]\centering\small
\begin{tabular}{@{}lp{0.72\textwidth}@{}}
\toprule
Setting & Output\\
\midrule
full marks & \ar{إِنْ ذِسْ بَارْتْ أُفْ هِيلِنْ رُولِنْزْ سَايْكُوسِينِمَا، شِي دِلْفْزْ إِنْتُو ذَا نَيْتْشَر أُفْ لَاكْ.}\\[3pt]
computed, full & \ar{إِنْ ذِسْ بَارْتْ أُفْ هِيلِنْ رُولِنْزْ سَايْكُوسِينِمَا، شِي دِلْفْزْ إِنْتُو ذَا نَيْتْشَرْ أُفْ لَاكْ.}\\[3pt]
medial sukoon & \ar{إِن ذِس بَارْت أُف هِيلِن رُولِنْز سَايْكُوسِينِمَا، شِي دِلْفْز إِنْتُو ذَا نَيْتْشَر أُف لَاك.}\\[3pt]
vowel marks only & \ar{إِن ذِس بَارت أُف هِيلِن رُولِنز سَايكُوسِينِمَا، شِي دِلفز إِنتُو ذَا نَيتشَر أُف لَاك.}\\[3pt]
unmarked & \ar{إن ذس بارت أف هيلن رولنز سايكوسينما، شي دلفز إنتو ذا نيتشر أف لاك.}\\[3pt]
v as feh, g as jeem & \ar{إِنْ ذِسْ بَارْتْ أُفْ هِيلِنْ رُولِنْزْ سَايْكُوسِينِمَا، شِي دِلْفْزْ إِنْتُو ذَا نَيْتْشَرْ أُفْ لَاكْ.}\\[3pt]
\bottomrule
\end{tabular}
\caption{One sentence in several settings. The second row is the computed form with the defaults.}\label{tab:demo}
\end{table}

\section{The density ladder}
The four densities are best seen on a single sentence. Each line contains the same letters; only the marks differ.

\begin{table}[ht]\centering\small
\begin{tabular}{@{}ll@{}}
\toprule
Density & Result\\
\midrule
full & \ar{ذَا لِتَرْزْ تْشَيْنْجْ، بَتْ ذَا وُرْدْزْ رِمَيْنْ.}\\[3pt]
medial & \ar{ذَا لِتَرْز تْشَيْنْج، بَت ذَا وُرْدْز رِمَيْن.}\\[3pt]
vowels & \ar{ذَا لِتَرز تشَينج، بَت ذَا وُردز رِمَين.}\\[3pt]
none & \ar{ذا لترز تشينج، بت ذا وردز رمين.}\\[3pt]
\bottomrule
\end{tabular}
\caption{One sentence at each density.}\label{tab:density}
\end{table}

\section{Worked passages}
The following passages are public-domain texts and short originals, set with the default options. Each block shows the Arabic-script form, a Latin respelling that is only a reading aid, and the English. They illustrate how the rules behave on running prose and not on isolated words, and they include some of the weaknesses that Chapter~\ref{ch:open} lists.

\subsection{Children's prose}
\begin{quote}\noindent \ar{أَلِسْ وَازْ بِغِنِنْغْ تُو غِتْ فِيرِي تَايَرْدْ أُفْ سِتِنْغْ بَايْ هِرْ سِسْتَرْ أُونْ ذَا بَانْغْكْ، أَنْدْ أُفْ هَافِنْغْ نُوثِنْغْ تُو دُو.}\\[2pt]
\ar{ألس واز بغننغ تو غت فيري تايرد أف ستنغ باي هر سستر أون ذا بانغك، أند أف هافنغ نوثنغ تو دو.}\\[1pt]
\textit{alus waaz bigining too get veree tayerd uv siting bay her sister aan dhu bangk und uv having nuthing too doo}\\[1pt]
Alice was beginning to get very tired of sitting by her sister on the bank, and of having nothing to do.
\end{quote}
\begin{quote}\noindent \ar{وُونْسْ أُورْ تْوَايْسْ شِي هَادْ بِيبْتْ إِنْتُو ذَا بُوكْ هِرْ سِسْتَرْ وَازْ رِيدِنْغْ، بَتْ إِتْ هَادْ نُو بِكْتْشَرْزْ أُورْ كُونْفَرْسَيْشَنْزْ إِنْ إِتْ.}\\[2pt]
\ar{وونس أور توايس شي هاد بيبت إنتو ذا بوك هر سستر واز ريدنغ، بت إت هاد نو بكتشرز أور كونفرسيشنز إن إت.}\\[1pt]
\textit{wuns awr tways shee had peept intoo dhu buk her sister waaz reeding but it had noh pikcherz awr kaanverseyshunz in it}\\[1pt]
Once or twice she had peeped into the book her sister was reading, but it had no pictures or conversations in it.
\end{quote}
\begin{quote}\noindent \ar{أَنْدْ وَتْ إِزْ ذَا يُوسْ أُفْ أَ بُوكْ، ثُوتْ أَلِسْ، وِثَاوْتْ بِكْتْشَرْزْ أُورْ كُونْفَرْسَيْشَنْزْ؟}\\[2pt]
\ar{أند وت إز ذا يوس أف أ بوك، ثوت ألس، وثاوت بكتشرز أور كونفرسيشنز؟}\\[1pt]
\textit{und wut iz dhu yoos uv u buk thawt alus withowt pikcherz awr kaanverseyshunz}\\[1pt]
And what is the use of a book, thought Alice, without pictures or conversations?
\end{quote}

\subsection{Scripture}
\begin{quote}\noindent \ar{إِنْ ذَا بِغِنِنْغْ غُودْ كْرَيْتَدْ ذَا هِيفِنْ أَنْدْ ذَا أَرْثْ.}\\[1pt]
\textit{in dhu bigining gaad kreeeytud dhu hevun und dhu erth}\\[1pt]
In the beginning God created the heaven and the earth.
\end{quote}
\begin{quote}\noindent \ar{أَنْدْ ذَا أَرْثْ وَازْ وِثَاوْتْ فُورْمْ، أَنْدْ فُويْدْ؛ أَنْدْ دَارْكْنِسْ وَازْ أَبُونْ ذَا فَيْسْ أُفْ ذَا دِيبْ.}\\[1pt]
\textit{und dhu erth waaz withowt fawrm und voyd und daarknus waaz upaan dhu feys uv dhu deep}\\[1pt]
And the earth was without form, and void; and darkness was upon the face of the deep.
\end{quote}
\begin{quote}\noindent \ar{أَنْدْ غُودْ سِدْ، لِتْ ذِرْ بِي لَايْتْ: أَنْدْ ذِرْ وَازْ لَايْتْ.}\\[1pt]
\textit{und gaad sed let dher bee layt und dher waaz layt}\\[1pt]
And God said, Let there be light: and there was light.
\end{quote}

\subsection{Verse}
\begin{quote}\noindent \ar{شَالْ آيْ كُمْبِرْ ذِي تُو أَ سَمَرْزْ دَيْ؟}\\[1pt]
\textit{shal ay kumper dhee too u sumerz dey}\\[1pt]
Shall I compare thee to a summer's day?
\end{quote}
\begin{quote}\noindent \ar{ذَاوْ آرْتْ مُورْ لَفْلِي أَنْدْ مُورْ تِمْبْرَتْ.}\\[1pt]
\textit{dhow aart mawr luvlee und mawr temprut}\\[1pt]
Thou art more lovely and more temperate.
\end{quote}
\begin{quote}\noindent \ar{رُوفْ وِنْدْزْ دُو شَيْكْ ذَا دَارْلِنْغْ بَدْزْ أُفْ مَيْ، أَنْدْ سَمَرْزْ لِيسْ هَاثْ أُولْ تُو شُورْتْ أَ دَيْتْ.}\\[1pt]
\textit{ruf windz doo sheyk dhu daarling budz uv mey und sumerz lees hath awl too shawrt u deyt}\\[1pt]
Rough winds do shake the darling buds of May, and summer's lease hath all too short a date.
\end{quote}
\begin{quote}\noindent \ar{وُونْسْ أَبُونْ أَ مِدْنَايْتْ دْرِيرِي، وَايْلْ آيْ بَانْدَرْدْ، وِيكْ أَنْدْ وِيرِي،}\\[1pt]
\textit{wuns upaan u midnayt driree wayl ay paanderd week und wiree}\\[1pt]
Once upon a midnight dreary, while I pondered, weak and weary,
\end{quote}
\begin{quote}\noindent \ar{أُوفَرْ مِنِي أَ كْوَيْنْتْ أَنْدْ كْيُرِيَسْ فُولْيُومْ أُفْ فَرْغُوتِنْ لُورْ.}\\[1pt]
\textit{ohver menee u kweynt und kyureeus vaalyoom uv fergaatun lawr}\\[1pt]
Over many a quaint and curious volume of forgotten lore.
\end{quote}

\subsection{A long sentence}
\begin{quote}\noindent \ar{وِنْ إِنْ ذَا كُورْسْ أُفْ هْيُومَنْ إِفِنْتْسْ، إِتْ بِكَمْزْ نِيسِسِرِي فُورْ وُونْ بِيبَلْ تُو دِزُولْفْ ذَا بُلِيتِكَلْ بَانْدْزْ وِتْشْ هَافْ كُنِكْتِدْ ذِمْ وِذْ أَنُوذَرْ، أَنْدْ تُو أَسُومْ أَمُونْغْ ذَا بَاوَرْزْ أُفْ ذَا أَرْثْ، ذَا سِيبَرَيْتْ أَنْدْ إِيكْوَلْ سْتَيْشَنْ تُو وِتْشْ ذَا لُوزْ أُفْ نَيْتْشَرْ أَنْدْ أُفْ نَيْتْشَرْزْ غُودْ إِنْتَايْتَلْ ذِمْ، أَ دِيسِنْتْ رِسْبِكْتْ تُو ذَا أُبِينْيَنْزْ أُفْ مَانْكَايْنْدْ رِكْوَايَرْزْ ذَتْ ذَيْ شُودْ دِكْلِرْ ذَا كَازِزْ وِتْشْ إِمْبِلْ ذِمْ تُو ذَا سِيبَرَيْشَنْ.}\\[1pt]
\textit{wen in dhu kawrs uv hyoomun ivents it bikumz nesuseree fawr wun peepul too dizaalv dhu pulitukul bandz wich hav kunektid dhem widh unudher und too usoom umung dhu powerz uv dhu erth dhu sepereyt und eekwul steyshun too wich dhu lawz uv neycher und uv neycherz gaad entaytul dhem u deesunt rispekt too dhu upinyunz uv mankaynd reekwayerz dhat dhey shud dikler dhu kaazuz wich impel dhem too dhu sepereyshun}\\[1pt]
When in the Course of human events, it becomes necessary for one people to dissolve the political bands which have connected them with another, and to assume among the powers of the earth, the separate and equal station to which the Laws of Nature and of Nature's God entitle them, a decent respect to the opinions of mankind requires that they should declare the causes which impel them to the separation.
\end{quote}

\subsection{Narrative}
\begin{quote}\noindent \ar{ذَا تْرَافِلَرْ كَيْمْ دَاوْنْ فْرُومْ ذَا سْتَارْ سِرَسْ تُو آوَرْ لِتَلْ أَنْتْ هِلْ، أَنْدْ مِزَرْدْ ذَا أَرْثْ وِذْ هِزْ فِنْغَرْزْ.}\\[1pt]
\textit{dhu travuler keym down frum dhu staar sirius too ower litul ant hil und mezherd dhu erth widh hiz finggerz}\\[1pt]
The traveller came down from the star Sirius to our little ant hill, and measured the earth with his fingers.
\end{quote}
\begin{quote}\noindent \ar{هِي فَاوْنْدْ ذَا بِيبَلْ سْمُولْ، بَتْ هِي فَاوْنْدْ ذِمْ كْيُرِيَسْ، أَنْدْ هِي أَسْكْتْ ذِمْ مِنِي كْوِسْتْشَنْزْ أَبَاوْتْ ذِرْ لِفْزْ.}\\[1pt]
\textit{hee fownd dhu peepul smawl but hee fownd dhem kyureeus und hee askt dhem menee kweschunz ubowt dher livz}\\[1pt]
He found the people small, but he found them curious, and he asked them many questions about their lives.
\end{quote}

\section{Names, dates and numbers}
Proper names are looked up like other words, and the fallback handles names the dictionary lacks. Numbers can be written in Arabic-Indic digits, in Latin digits or spoken, depending on the \texttt{digits} option. The spoken form is useful for text meant to be read aloud.

\begin{quote}\noindent \ar{هِيلِنْ رُولِنْزْ لِفْزْ إِنْ هَالِفَاكْسْ، أَنْدْ هِرْ بْرُوذَرْ فِكْتَرْ لِفْزْ إِنْ فَنْكُوفَرْ.}\\[1pt]
\textit{helun raalinz livz in halifaks und her brudher vikter livz in vankoover}\\[1pt]
Helen Rollins lives in Halifax, and her brother Victor lives in Vancouver.
\end{quote}
\begin{quote}\noindent \ar{تُوزْدِي، ذَا فِفْثْ أُفْ أُوكْتُوبَرْ، أَتْ سِيفِنْ ثِرْدِي، ذَا تْرَيْنْ فْرُومْ تَرُونْتُو أَرَايْفْزْ.}\\[1pt]
\textit{toozdee dhu fifth uv aaktohber at sevun therdee dhu treyn frum teraantoh erayvz}\\[1pt]
Tuesday, the fifth of October, at seven thirty, the train from Toronto arrives.
\end{quote}
\begin{quote}\noindent \ar{ذِرْ أَرْ ٢٥ وُرْدْزْ أُونْ ذِسْ بَيْجْ أَنْدْ ٣ أُفْ ذِمْ أَرْ نَمْبَرْزْ.}\\[1pt]
\textit{dher aar werdz aan dhis peyj und uv dhem aar numberz}\\[1pt]
There are 25 words on this page and 3 of them are numbers.
\end{quote}
\begin{quote}\noindent \ar{إِنْ ٢٠٢٦ إِتْ كُوسْتْ ١,٢٥٠ دُولَرْزْ.}\\[1pt]
\textit{in it kaast daalerz}\\[1pt]
In 2026 it cost 1,250 dollars.
\end{quote}

\section{Reading the results}
Three habits of the output are worth noting. Common words always take their fixed form, so \emph{the}, \emph{of}, \emph{and} and \emph{is} look the same everywhere. Stressed vowels are usually marked more heavily than unstressed ones, which gives longer words a visible rhythm. And function words with weak forms in speech, such as \emph{to} and \emph{a}, keep their fixed spellings and are not reduced.

\chapter{The Reverse Reader}\label{ch:reader}

\section{Why the notation is many to one}
Several English spellings share one sound, and several English sounds share one letter. The notation therefore cannot return a unique English word without help. \emph{knight} and \emph{night} are written identically, and so are \emph{write} and \emph{right}. Unstressed vowels that are optional marks add further collisions, as do the shared letters \ar{ب} (for $b$ and $p$, and for $v$ under the \texttt{b} setting) and \ar{ج} or \ar{غ} (for sounds merged with $j$).

\section{Reader design}
The reader indexes every dictionary word by its Arabic skeleton, the letters without marks, and by a coarser key without vowel letters. For an input word it returns candidates in tiers: attested forms first, then skeleton matches ranked by corpus frequency with a preference for those whose marks agree with the input, then matches on the coarse key. It reports a ranked list and never a single guess.

Table~\ref{tab:rt} shows pairs that the notation cannot separate. The original word is among the candidates but not always first.

\begin{table}[ht]\centering\small
\begin{tabular}{@{}llll@{}}
\toprule
English & Spelling & Top candidates & Rank of original\\
\midrule
knight & \ar{نَايْتْ} & night, knight, nite & 2\\
planet & \ar{بْلَانِتْ} & plant, planet, plante, pylant & 2\\
their & \ar{ذِرْ} & their, there, thur & 1\\
write & \ar{رَايْتْ} & right, write, wright, riot & 2\\
\bottomrule
\end{tabular}
\caption{Words whose spelling is shared with another English word.}\label{tab:rt}
\end{table}

\section{Measured ambiguity}
Two measures describe the notation, both over the \NFreq{} most frequent English words that the dictionary covers.

\begin{itemize}
\item \textbf{Collisions.} \PctSharedFull{}\% of these words share their fully marked form with at least one other word, and \PctSharedSkel{}\% share their unmarked form. Table~\ref{tab:coll} shows the largest classes.
\item \textbf{Reading back.} For the \RtWords{} most frequent words, the reader returns the original as its first candidate \RtTopOne{}\% of the time and among its first five \RtTopFive{}\%. The mean candidate list has \RtMeanCand{} entries. This figure is optimistic by construction: the test words are frequent, frequency ranks the candidates, and the reader uses the same dictionary as the writer.
\end{itemize}

\begin{table}[ht]\centering\small
\begin{tabular}{@{}lp{0.6\textwidth}@{}}
\toprule
Unmarked form & Words sharing it\\
\midrule
\ar{بت} & but, put, bit, bet, pet, butt, pit\\
\ar{بول} & ball, paul, pool, bowl, poll, pole\\
\ar{بي} & be, pay, b, p, bay\\
\ar{بن} & been, ben, pen, pin, bin\\
\ar{سي} & see, say, c, sea, se\\
\ar{أن} & an, un, ann, anne\\
\bottomrule
\end{tabular}
\caption{Largest classes of words that share an unmarked spelling.}\label{tab:coll}
\end{table}

\section{Consequence}
Exact recovery of English spelling from the Arabic form is not available. A phonetic reading is always available. A lossless round trip would require extra information: surrounding context, a language model over the candidates, or optional disambiguation marks. The tools therefore separate recovering a pronunciation from recovering a spelling, and they report ambiguity in place of hiding it.

\section{A reading in detail}
The reader treats each token separately. The example below passes the earlier sentence back through it, with the specimen forms switched off. Fixed spellings of frequent words are always available, because they are part of the notation.

\begin{table}[ht]\centering\small
\begin{tabular}{@{}lll@{}}
\toprule
Written & Top candidates & Reader tier\\
\midrule
\ar{ذَا} & the, with, they, though & attested\\
\ar{لِتَرْزْ} & letters, walters & skeleton\\
\ar{تْشَيْنْجْ،} & change & skeleton\\
\ar{بَتْ} & but, butt, put, bit & skeleton\\
\ar{ذَا} & the, with, they, though & attested\\
\ar{وُرْدْزْ} & words, awards, roads, yards & skeleton\\
\ar{رِمَيْنْ.} & remain, roman, romania, romney & skeleton\\
\bottomrule
\end{tabular}
\caption{Reader output for one sentence: the first four candidates for each token.}\label{tab:readerdemo}
\end{table}

In most rows the intended word is first. The token \ar{إِزْ} is shared by \emph{is} and \emph{as}, and the reader cannot choose between them without context. Such pairs are the main cost of the compact notation.

\section{Using context}
The reader treats words independently. Choosing among candidates with a language model over the whole sentence is the natural next step, because each candidate list is short and most ambiguities resolve easily from context. The data structures already return ranked lists, so a context model can be added as a reranking step without changing the index.

\section{Tiers and their meaning}
The tier attached to each candidate says how closely it matches. An attested or fixed form is an exact recovery. A skeleton match means the letters agree and the marks are compatible. A consonant match means only the consonant skeleton agrees, and it is offered when the other tiers have little to say. A reader that shows the tier lets the user judge how much weight a candidate deserves.

\chapter{The Linter}

\section{Purpose}
The specimens drift. The Atlas text uses Persian letters, the Micromegas text uses one of them, a web export replaced some letters with the Unicode replacement character, and typed passages occasionally place two vowel marks on one letter. A notation that depends on a small inventory needs a mechanical check that the inventory is respected. The linter in \texttt{arabglish.lint} provides that check.

\section{Checks}
The linter reports each problem with its position, the offending character, a code and a message. The codes are as follows.

\begin{description}
\item[extended-letter] A letter outside the standard Arabic block, such as \ar{پ}, \ar{چ}, \ar{ژ}, \ar{گ} or \ar{ڤ}. A suggested replacement is attached.
\item[corrupt] The replacement character U+FFFD. The original letter is lost and cannot be repaired automatically.
\item[orphan-mark] A vowel mark or sukoon with no base letter.
\item[double-vowel] Two vowel marks on one letter.
\item[sukoon-and-vowel] A sukoon and a vowel mark on one letter.
\item[shadda] A shadda. The system does not mark English doubling, so a shadda is an error.
\item[tatweel] The elongation character, which is not part of the system.
\item[unknown-arabic] Any other Arabic-block character outside the inventory.
\end{description}

Punctuation, Latin text, digits and the Arabic comma, semicolon and question mark are accepted.

\section{Repair}
Only the deterministic part of the repair is automated. The function \texttt{fix} rewrites extended letters to the core inventory: \ar{پ} becomes \ar{ب}, \ar{ڤ} becomes \ar{ف}, \ar{گ} becomes \ar{غ}, \ar{چ} becomes \ar{تْش}, \ar{ژ} becomes \ar{ز}, and the Persian forms of kaf and yeh become the Arabic ones. This follows the decisions in Chapter~\ref{ch:evidence}. If a different decision is adopted for $v$ or $g$, the replacement table is the one place to change. Corrupt characters, orphan marks and double marks are reported but left for the author, because several repairs are possible.

\section{Results on the specimens}
Table~\ref{tab:inventory} lists the issue counts. The page transcript and typed passage are clean. The Atlas text has \AtlasExtended{} issues, nearly all extended letters. Running \texttt{fix} on it produces text that passes the linter, which is a convenient way to bring early material into the current inventory.

\section{A worked example}
The line below is a deliberately damaged piece of text. It contains four extended letters, a stray shadda and an orphan mark. The linter lists every problem with its position, and \texttt{fix} repairs the part that has a single correct repair.

\begin{table}[ht]\centering\small
\begin{tabular}{@{}lll@{}}
\toprule
Character & Code & Suggested repair\\
\midrule
\ar{گ} & extended-letter & \ar{غ}\\
\ar{ک} & extended-letter & \ar{ك}\\
\ar{ڤ} & extended-letter & \ar{ف}\\
\ar{چ} & extended-letter & \ar{تْش}\\
\ar{پ} & extended-letter & \ar{ب}\\
\ar{ّ} & shadda & --\\
\ar{ّ} & orphan-mark & --\\
\bottomrule
\end{tabular}
\caption{Linter output for a damaged line.}\label{tab:lintdemo}
\end{table}

Before and after the deterministic repair:
\begin{quote}\noindent
\ar{ذَا گُود بُوک إِزْ ڤَرِي چِيپْ ّ} \quad $\rightarrow$ \quad \ar{ذَا غُود بُوك إِزْ فَرِي تْشِيبْ ّ}
\end{quote}
The shadda is not repaired because there is no single correct reading of it. It may be an error or a deliberate marking of a doubled consonant, and the author decides which.

\section{Using the linter in a workflow}
The linter returns a nonzero status when issues are found, so it can run in a script before typesetting or before a text is added to the specimen set. A typical sequence is to repair the extended letters automatically, rerun the linter, and read the remaining items by hand. Because the linter reports positions, an editor can jump to each item.

\section{What the linter does not check}
The linter checks the inventory and the order of marks. It does not check that the spelling is a good rendering of any English word, because the notation has no single correct spelling for many words. Checking against a pronunciation is the job of the forward transliterator, and the evaluation in Chapter~\ref{ch:eval} compares its output with attested forms.

\chapter{Evaluation}\label{ch:eval}

\section{What can be measured}
The only ground truth available is the set of \NAttested{} word forms written by the author in the typed passage and the handwritten page. The computed form of each word is compared with the attested form at two levels: exact agreement of letters and marks, and agreement of letters once marks are removed.

\section{Results}
Of \NAttested{} attested words, \NExact{} (\PctExact\%) are reproduced exactly and \NLetters{} (\PctLetters\%) agree at the level of letters. Ten words differ. Table~\ref{tab:differs} lists them.

The differences fall into five groups. A lexical alef appears in forms such as \emph{capitalism}, where the author writes a long vowel the dictionary reading does not require. The choice for $v$ differs in \emph{individual} and \emph{universal}, where the typed passage uses \ar{ب} and the default is \ar{ف}. Some alignments of spelling to sound fail on longer words such as \emph{individual} and \emph{ideology}. The seat letter for an initial vowel differs in \emph{a}. Finally, the digraph \emph{ch} in \emph{nature} is written \ar{تْش} on the page and \ar{ش} in the typed passage, so one of the two attested forms must disagree with any single rule.

\section{Interpretation}
The rules of Chapter~\ref{ch:vowels} were fitted to these same specimens. The agreement is therefore a description of how well the rules summarise the evidence and is not a validation on unseen text. Held-out specimens are the main need; Chapter~\ref{ch:open} returns to this.

\section{Regression ratchet}
The test suite fixes lower bounds on the two agreement figures (0.77 for letters and 0.72 for exact forms). A change to the rules that lowers either figure fails the tests. A change that raises them is accepted and the bound can be moved up. This keeps later work from silently undoing earlier fits.

\section{Evaluating the reader}
The reader is evaluated separately in Chapter~\ref{ch:reader}. It is measured with the attested tier switched off, so that specimen words do not recognise themselves, and on frequent words only.

\section{Sensitivity to the open switches}
The two main open switches change few attested words, so the agreement figures barely move with them. Table~\ref{tab:sens} shows all four combinations of the settings for $v$ and hard $g$.

\begin{table}[ht]\centering\small
\begin{tabular}{@{}llrr@{}}
\toprule
$v$ & hard $g$ & Exact & Letters\\
\midrule
\ar{ف} & \ar{غ} & 33 & 35\\
\ar{ف} & \ar{ج} & 33 & 35\\
\ar{ب} & \ar{غ} & 33 & 35\\
\ar{ب} & \ar{ج} & 33 & 35\\
\bottomrule
\end{tabular}
\caption{Agreement with the \NAttested{} attested forms under each setting of $v$ and hard $g$.}\label{tab:sens}
\end{table}

The equal numbers do not mean the settings are equivalent. The attested set contains only two $v$ words with disagreeing letters, and no word with a hard $g$, so the specimens cannot rank the options. This is the reason both stay as switches.

\section{Agreement on newly composed examples}
After the main rules were fixed, three further sentences were composed in the notation and reviewed by the author, who accepted them with one correction (\emph{book}). They were then compared with the toolkit word by word. Of \NApprovedWords{} words, \NApprovedSame{} are reproduced exactly. Table~\ref{tab:appdiff} lists the others.

\begin{table}[ht]\centering\small
\begin{tabular}{@{}lll@{}}
\toprule
English & Approved & Computed\\
\midrule
words & \ar{وَرْدْزْ} & \ar{وُرْدْزْ}\\
recognise & \ar{رِكِغْنَايْزْ} & \ar{رِيكُغْنَايْزْ}\\
familiar & \ar{فَمِلْيَرْ} & \ar{فَمِيلْيَرْ}\\
\bottomrule
\end{tabular}
\caption{Words in the approved sentences that the toolkit writes differently.}\label{tab:appdiff}
\end{table}

These sentences were composed after the rules and not before, but they were composed by the same system of reasoning and are not independent evidence. Their value is as a check that the rules produce forms the author finds acceptable on unseen words. One difference, \emph{words}, concerns the vowel before a cluster, and two more concern the open-syllable yeh in \emph{recognise} and \emph{familiar}. Each is a candidate for a rule change once the author has decided which form is wanted.

\section{Error analysis}
The remaining disagreements with the attested forms are of three kinds. Some are lexical, such as the alef in \emph{capitalism}, and no rule predicts them. Some come from alignment, where the number of vowel letter groups in the spelling does not match the number of vowel sounds. And some are choices that the specimens contradict among themselves, such as \emph{nature}. Only the second kind can be reduced by better software, and the first and third require decisions or more data.

\section{A proper held-out test}
A real test needs words the rules have not seen. The simplest design is to give the author a list of unseen words, have each written by hand in the notation without the tool, and compare. Fifty words chosen across frequency bands, with a few names and a few long technical words, would be enough to show which rules generalise and which were overfitted. The scripts already accept such a file and print the same tables as for the present specimens.

\chapter{Specimen Glossary}
This chapter applies the forward transliterator to the opening of the typed passage. Each entry has three lines: the Arabic-script form with the default settings, a Latin respelling generated from the Arabic form by the guide tool, and the English original. The respelling is a reading aid for the author and for readers who do not read Arabic script. It is not part of the notation.

\begin{quote}\noindent \ar{ذَا رِبْرِشَنْ أُفْ لَاكْ أَنْدْ إِتْسْ كَلْتْشَرَلْ سِمْبْتُمْزْ}\\[1pt]
\textit{dhu reepreshun uv lak und its kulcherul simptumz}\\[1pt]
The repression of lack and its cultural symptoms.
\end{quote}
\begin{quote}\noindent \ar{إِنْ ذِسْ بَارْتْ أُفْ هِيلِنْ رُولِنْزْ سَايْكُوسِينِمَا}\\[1pt]
\textit{in dhis paart uv helun raalinz saykohsinumu}\\[1pt]
In this part of Helen Rollins' psychocinema,
\end{quote}
\begin{quote}\noindent \ar{شِي دِلْفْزْ إِنْتُو ذَا نَيْتْشَر أُفْ لَاكْ}\\[1pt]
\textit{shee delvz intoo dhu neycher uv lak}\\[1pt]
she delves into the nature of lack.
\end{quote}

\section{Ambiguity and context}
Some distinctions disappear in the notation. \emph{Is} and \emph{as} share a written form, and so do \emph{night} and \emph{knight}. Taken one at a time, such collisions show where the system's limits lie, and Chapter~\ref{ch:reader} measures them. Their practical weight is a separate question, because readers rarely meet words in isolation. In a sentence, the intended word is usually clear from its neighbours, as it is for homophones in spoken English, and a collision matters only where it also causes confusion in continuous prose. The lists below are therefore best read as evidence about the notation's resolution, and not as a prediction that each pair will cause trouble in use.

\section{The reference vocabulary}
The tables below group words by the feature they illustrate. Every form is generated by the toolkit with the default options, and every table is regenerated whenever the rules change, so the book cannot drift from the code.

\subsection{Doubled o and related spellings}
\begin{tabular}{@{}lll@{}}
\toprule
English & Arabic script & Respelling\\
\midrule
book & \ar{بُوكْ} & \textit{buk}\\
look & \ar{لُوكْ} & \textit{luk}\\
good & \ar{غُودْ} & \textit{gud}\\
foot & \ar{فُوتْ} & \textit{fut}\\
wood & \ar{وُودْ} & \textit{wud}\\
food & \ar{فُودْ} & \textit{food}\\
moon & \ar{مُونْ} & \textit{moon}\\
school & \ar{سْكُولْ} & \textit{skool}\\
too & \ar{تُو} & \textit{too}\\
two & \ar{تُو} & \textit{too}\\
blue & \ar{بْلُو} & \textit{bloo}\\
true & \ar{تْرُو} & \textit{troo}\\
\bottomrule
\end{tabular}

\subsection{ee and ea}
\begin{tabular}{@{}lll@{}}
\toprule
English & Arabic script & Respelling\\
\midrule
see & \ar{سِي} & \textit{see}\\
tree & \ar{تْرِي} & \textit{tree}\\
green & \ar{غْرِينْ} & \textit{green}\\
sleep & \ar{سْلِيبْ} & \textit{sleep}\\
head & \ar{هِدْ} & \textit{hed}\\
bread & \ar{بْرِدْ} & \textit{bred}\\
read & \ar{رِدْ} & \textit{red}\\
meet & \ar{مِيتْ} & \textit{meet}\\
meat & \ar{مِيتْ} & \textit{meet}\\
dream & \ar{دْرِيمْ} & \textit{dreem}\\
dead & \ar{دِدْ} & \textit{ded}\\
\bottomrule
\end{tabular}

\subsection{Glides and diphthongs}
\begin{tabular}{@{}lll@{}}
\toprule
English & Arabic script & Respelling\\
\midrule
rain & \ar{رَيْنْ} & \textit{reyn}\\
day & \ar{دَيْ} & \textit{dey}\\
light & \ar{لَايْتْ} & \textit{layt}\\
night & \ar{نَايْتْ} & \textit{nayt}\\
voice & \ar{فُويْسْ} & \textit{voys}\\
flower & \ar{فْلَاوَرْ} & \textit{flower}\\
hour & \ar{آوَرْ} & \textit{ower}\\
boy & \ar{بُويْ} & \textit{boy}\\
toy & \ar{تُويْ} & \textit{toy}\\
noise & \ar{نُويْزْ} & \textit{noyz}\\
now & \ar{نَاوْ} & \textit{now}\\
how & \ar{هَاوْ} & \textit{how}\\
\bottomrule
\end{tabular}

\subsection{Short vowels}
\begin{tabular}{@{}lll@{}}
\toprule
English & Arabic script & Respelling\\
\midrule
cat & \ar{كَاتْ} & \textit{kat}\\
bed & \ar{بِدْ} & \textit{bed}\\
sit & \ar{سِتْ} & \textit{sit}\\
hot & \ar{هُوتْ} & \textit{haat}\\
cut & \ar{كَتْ} & \textit{kut}\\
put & \ar{بُتْ} & \textit{put}\\
pen & \ar{بِنْ} & \textit{pen}\\
map & \ar{مَابْ} & \textit{map}\\
top & \ar{تُوبْ} & \textit{taap}\\
sun & \ar{سَنْ} & \textit{sun}\\
\bottomrule
\end{tabular}

\subsection{Silent e}
\begin{tabular}{@{}lll@{}}
\toprule
English & Arabic script & Respelling\\
\midrule
make & \ar{مَيْكْ} & \textit{meyk}\\
time & \ar{تَايْمْ} & \textit{taym}\\
home & \ar{هُومْ} & \textit{hohm}\\
use & \ar{يُوسْ} & \textit{yoos}\\
cake & \ar{كَيْكْ} & \textit{keyk}\\
like & \ar{لَايْكْ} & \textit{layk}\\
phone & \ar{فُونْ} & \textit{fohn}\\
tune & \ar{تُونْ} & \textit{toon}\\
\bottomrule
\end{tabular}

\subsection{R-coloured vowels}
\begin{tabular}{@{}lll@{}}
\toprule
English & Arabic script & Respelling\\
\midrule
car & \ar{كَارْ} & \textit{kaar}\\
far & \ar{فَارْ} & \textit{faar}\\
her & \ar{هِرْ} & \textit{her}\\
bird & \ar{بِرْدْ} & \textit{berd}\\
turn & \ar{تَرْنْ} & \textit{tern}\\
word & \ar{وُرْدْ} & \textit{werd}\\
door & \ar{دُورْ} & \textit{dawr}\\
more & \ar{مُورْ} & \textit{mawr}\\
fire & \ar{فَايَرْ} & \textit{fayer}\\
tire & \ar{تَايَرْ} & \textit{tayer}\\
\bottomrule
\end{tabular}

\subsection{Consonants}
\begin{tabular}{@{}lll@{}}
\toprule
English & Arabic script & Respelling\\
\midrule
thing & \ar{ثِنْغْ} & \textit{thing}\\
think & \ar{ثِنْغْكْ} & \textit{thingk}\\
this & \ar{ذِسْ} & \textit{dhis}\\
that & \ar{ذَتْ} & \textit{dhat}\\
church & \ar{تْشَرْتْشْ} & \textit{cherch}\\
change & \ar{تْشَيْنْجْ} & \textit{cheynj}\\
judge & \ar{جَجْ} & \textit{juj}\\
measure & \ar{مِيزَرْ} & \textit{mezher}\\
vision & \ar{فِيزَنْ} & \textit{vizhun}\\
singer & \ar{سِنْغَرْ} & \textit{singer}\\
finger & \ar{فِنْغَرْ} & \textit{fingger}\\
garden & \ar{غَارْدِنْ} & \textit{gaardun}\\
\bottomrule
\end{tabular}

\subsection{Clusters}
\begin{tabular}{@{}lll@{}}
\toprule
English & Arabic script & Respelling\\
\midrule
string & \ar{سْتْرِنْغْ} & \textit{string}\\
spring & \ar{سْبْرِنْغْ} & \textit{spring}\\
splash & \ar{سْبْلَاشْ} & \textit{splash}\\
street & \ar{سْتْرِيتْ} & \textit{street}\\
strength & \ar{سْتْرِنْغْكْثْ} & \textit{strengkth}\\
twelfth & \ar{تْوِلْفْثْ} & \textit{twelfth}\\
asked & \ar{أَسْكْتْ} & \textit{askt}\\
texts & \ar{تِكْسْتْسْ} & \textit{teksts}\\
\bottomrule
\end{tabular}

\section{Teaching sequence}
The materials in this chapter suit a staged course in which the supporting lines are withdrawn as the reader gains fluency.
\begin{enumerate}
\item \textbf{Paired reading.} Begin with the Arabic-script English beside the familiar English word, so that each new spelling is tied to a word the reader already knows.
\item \textbf{A small recurring vocabulary.} Practise a few dozen frequent words (\emph{the}, \emph{of}, \emph{and}, \emph{is}, \emph{in}, \emph{to}) until their written forms are recognised at once. Their fixed spellings make this easy.
\item \textbf{Short passages with the English covered.} Read short passages with the English line concealed, uncovering it only to check.
\item \textbf{Withdraw the respelling.} Drop the Latin respelling next, since it is an approximate aid. It introduces spelling conventions of its own (the \emph{ee} of \emph{reeding}, the \emph{u} of \emph{uv}), and a reader who leans on it too long learns those conventions instead of the notation.
\item \textbf{Lighter marks.} Finally, move from full marks to medial, vowel-only and unmarked text, using the density settings.
\end{enumerate}

\section{Reading practice}
The following passages start with the simplest sentences and become longer. Readers learning the notation can cover the English and the Latin respelling and read the Arabic-script line aloud.

\begin{quote}\noindent \ar{ذَا بُوكْ إِزْ رِتِنْ إِنْ إِنْغْلِشْ.}\\[2pt]
\ar{ذا بوك إز رتن إن إنغلش.}\\[1pt]
\textit{dhu buk iz ritun in ingglish}\\[1pt]
The book is written in English.
\end{quote}
\begin{quote}\noindent \ar{ذَا لِتَرْزْ تْشَيْنْجْ، بَتْ ذَا وُرْدْزْ رِمَيْنْ.}\\[2pt]
\ar{ذا لترز تشينج، بت ذا وردز رمين.}\\[1pt]
\textit{dhu leterz cheynj but dhu werdz rimeyn}\\[1pt]
The letters change, but the words remain.
\end{quote}
\begin{quote}\noindent \ar{أَ رِيدَرْ لَرْنْزْ ذَا سْكْرِبْتْ أَنْدْ بِغِنْزْ تُو رِيكُغْنَايْزْ أَ فَمِيلْيَرْ فُويْسْ.}\\[2pt]
\ar{أ ريدر لرنز ذا سكربت أند بغنز تو ريكغنايز أ فميلير فويس.}\\[1pt]
\textit{u reeder lernz dhu skript und biginz too rekugnayz u fumilyer voys}\\[1pt]
A reader learns the script and begins to recognise a familiar voice.
\end{quote}
\begin{quote}\noindent \ar{أَلِسْ وَازْ بِغِنِنْغْ تُو غِتْ فِيرِي تَايَرْدْ أُفْ سِتِنْغْ بَايْ هِرْ سِسْتَرْ أُونْ ذَا بَانْغْكْ، أَنْدْ أُفْ هَافِنْغْ نُوثِنْغْ تُو دُو.}\\[2pt]
\ar{ألس واز بغننغ تو غت فيري تايرد أف ستنغ باي هر سستر أون ذا بانغك، أند أف هافنغ نوثنغ تو دو.}\\[1pt]
\textit{alus waaz bigining too get veree tayerd uv siting bay her sister aan dhu bangk und uv having nuthing too doo}\\[1pt]
Alice was beginning to get very tired of sitting by her sister on the bank, and of having nothing to do.
\end{quote}
\begin{quote}\noindent \ar{وُونْسْ أُورْ تْوَايْسْ شِي هَادْ بِيبْتْ إِنْتُو ذَا بُوكْ هِرْ سِسْتَرْ وَازْ رِيدِنْغْ، بَتْ إِتْ هَادْ نُو بِكْتْشَرْزْ أُورْ كُونْفَرْسَيْشَنْزْ إِنْ إِتْ.}\\[2pt]
\ar{وونس أور توايس شي هاد بيبت إنتو ذا بوك هر سستر واز ريدنغ، بت إت هاد نو بكتشرز أور كونفرسيشنز إن إت.}\\[1pt]
\textit{wuns awr tways shee had peept intoo dhu buk her sister waaz reeding but it had noh pikcherz awr kaanverseyshunz in it}\\[1pt]
Once or twice she had peeped into the book her sister was reading, but it had no pictures or conversations in it.
\end{quote}
\begin{quote}\noindent \ar{أَنْدْ وَتْ إِزْ ذَا يُوسْ أُفْ أَ بُوكْ، ثُوتْ أَلِسْ، وِثَاوْتْ بِكْتْشَرْزْ أُورْ كُونْفَرْسَيْشَنْزْ؟}\\[2pt]
\ar{أند وت إز ذا يوس أف أ بوك، ثوت ألس، وثاوت بكتشرز أور كونفرسيشنز؟}\\[1pt]
\textit{und wut iz dhu yoos uv u buk thawt alus withowt pikcherz awr kaanverseyshunz}\\[1pt]
And what is the use of a book, thought Alice, without pictures or conversations?
\end{quote}
\begin{quote}\noindent \ar{إِنْ ذَا بِغِنِنْغْ غُودْ كْرَيْتَدْ ذَا هِيفِنْ أَنْدْ ذَا أَرْثْ.}\\[1pt]
\textit{in dhu bigining gaad kreeeytud dhu hevun und dhu erth}\\[1pt]
In the beginning God created the heaven and the earth.
\end{quote}
\begin{quote}\noindent \ar{أَنْدْ ذَا أَرْثْ وَازْ وِثَاوْتْ فُورْمْ، أَنْدْ فُويْدْ؛ أَنْدْ دَارْكْنِسْ وَازْ أَبُونْ ذَا فَيْسْ أُفْ ذَا دِيبْ.}\\[1pt]
\textit{und dhu erth waaz withowt fawrm und voyd und daarknus waaz upaan dhu feys uv dhu deep}\\[1pt]
And the earth was without form, and void; and darkness was upon the face of the deep.
\end{quote}
\begin{quote}\noindent \ar{أَنْدْ غُودْ سِدْ، لِتْ ذِرْ بِي لَايْتْ: أَنْدْ ذِرْ وَازْ لَايْتْ.}\\[1pt]
\textit{und gaad sed let dher bee layt und dher waaz layt}\\[1pt]
And God said, Let there be light: and there was light.
\end{quote}
\begin{quote}\noindent \ar{ذَا تْرَافِلَرْ كَيْمْ دَاوْنْ فْرُومْ ذَا سْتَارْ سِرَسْ تُو آوَرْ لِتَلْ أَنْتْ هِلْ، أَنْدْ مِزَرْدْ ذَا أَرْثْ وِذْ هِزْ فِنْغَرْزْ.}\\[1pt]
\textit{dhu travuler keym down frum dhu staar sirius too ower litul ant hil und mezherd dhu erth widh hiz finggerz}\\[1pt]
The traveller came down from the star Sirius to our little ant hill, and measured the earth with his fingers.
\end{quote}
\begin{quote}\noindent \ar{هِي فَاوْنْدْ ذَا بِيبَلْ سْمُولْ، بَتْ هِي فَاوْنْدْ ذِمْ كْيُرِيَسْ، أَنْدْ هِي أَسْكْتْ ذِمْ مِنِي كْوِسْتْشَنْزْ أَبَاوْتْ ذِرْ لِفْزْ.}\\[1pt]
\textit{hee fownd dhu peepul smawl but hee fownd dhem kyureeus und hee askt dhem menee kweschunz ubowt dher livz}\\[1pt]
He found the people small, but he found them curious, and he asked them many questions about their lives.
\end{quote}

\chapter{Typesetting and Input}\label{ch:typesetting}

\section{Typesetting this book}
The book is set with XeLaTeX. The Arabic-script specimens are short runs inside English text, so the bidirectional machinery of a full Arabic setup is not needed. The \texttt{\textbackslash ar} macro starts a right-to-left run with the TeXXeT primitives (\texttt{\textbackslash beginR} and \texttt{\textbackslash endR}) and selects DejaVu Sans, which carries the Arabic block and positions fatha, damma, kasra and sukoon correctly. Packages that usually do this work, \texttt{polyglossia} and \texttt{bidi}, were not available in the build environment, and LuaLaTeX failed on font lookup. The TeXXeT route is small and has been checked by rendering pages and inspecting them.

Two cautions apply. A right-to-left run must not be broken across lines by hand, since the order of the pieces then depends on the line breaker. Letters and marks that arrive in the wrong order from a generator will render wrongly but silently, which is a reason to lint the text before typesetting it.

\section{Input}
The AutoHotkey script \texttt{toarabic.ahk} (Appendix~\ref{app:legacy}) is an input method: it maps one Latin key to one Arabic letter. That is useful for typing the notation by hand and unrelated to the sound key. A keyboard that matches the sound key would instead offer the consonant letters under the Latin keys that look most like the sounds (\texttt{b} and \texttt{p} on \ar{ب}, \texttt{f} on \ar{ف}, \texttt{g} on \ar{غ}) and the three vowel marks and sukoon on spare keys. Such a layout could be generated from the tables in \texttt{tables.py}. It is listed with the open work.

\section{Mark density in print}
Fully marked text is long and visually dense. The \texttt{medial} setting removes sukoon at the ends of words, \texttt{vowels} keeps vowel marks only, and \texttt{none} removes all marks. Because density is applied as a last filter, the same source text can be set at several densities without changing any letters. Table~\ref{tab:demo} shows the settings side by side.

\section{Fonts}
The notation needs a font that carries the basic Arabic block and positions the three vowel marks and sukoon cleanly above and below their base letters. DejaVu Sans does this reliably in XeLaTeX. Other fonts with full Arabic support, such as Amiri or Noto Naskh, can be substituted, and it is worth checking each by rendering a page, because some place the marks too close to the letters when several letters carry marks in a row.

\section{Copying and pasting}
Marks are separate code points that follow their base letter. Text copied from a web page or a word processor can lose or reorder them, and the linter's orphan-mark check catches the usual damage. Before a text is added to the specimens, it is worth running the linter once on the pasted result.

\section{Printing the book}
Because Arabic-script runs are short and set inline, the line breaker treats them as ordinary text, and long runs are better placed in their own paragraph or in a quotation block, as the worked passages in this book are. A run that must break across lines should be split between words, since splitting inside a word separates letters from their marks.

\section{Keyboard layout}
The AutoHotkey layout in the author's repository maps one Latin key to one Arabic letter, which is an input method. A layout designed from the sound key would put the letters under the Latin keys that look most like the sounds and place the marks on spare keys. A table-driven generator for such a layout would take its data from \texttt{tables.py}, so that a change of defaults changes the layout too. It is listed under open work, because it depends on the author's decision about how many marks should be typed by hand.

\section{Accessibility}
Screen readers read Arabic script with Arabic pronunciation rules, so notation text will be read as if it were Arabic. For an audience that needs speech, the Latin respelling or the English text should be offered next to it.

\chapter{Open Work}\label{ch:open}

\section{Pronunciation source}
The forward transliterator reads pronunciations from the CMU Pronouncing Dictionary, a General American dictionary with about 126\,000 entries, and falls back to a small spelling-based guesser for words it lacks. An earlier version of the system used a different pronunciation source, described as a kind of radio English, and the dictionary behind that version has not been identified. Which accent the notation should encode is an author decision. Candidates include another machine-readable dictionary (for example a British one with received pronunciation), or a table of per-word overrides for words where the author's own reading differs. The toolkit isolates the pronunciation source in one function, so changing it does not touch the sound key.

\section{The decisions still open}
Table~\ref{tab:decisions} lists the switches. The most consequential is the hard $g$ (\ar{غ} or \ar{ج}); $v$ is now \ar{ف} by the author's decision, with \ar{ب} available. The hard $g$ can be settled by a statement from the author or by a larger sample of attested forms.

\section{More specimens}
All fitted rules come from one typed passage and one handwritten page. Held-out hand-written pages, written without reference to the tool, would allow a real test of the rules in Chapter~\ref{ch:vowels}.

\section{Disambiguation}
Chapter~\ref{ch:reader} showed that the notation merges some words. Optional disambiguation marks, or a language model that ranks reader candidates in context, would reduce this. Whether the notation should carry such marks at all is a design choice that the author has not yet made.

\section{Further questions}
Whether English doubling should ever be marked with shadda, how to write stress, how to treat names from other languages, and whether the system should have dialect profiles are all left open. None is needed for the present tools.

\section{Priorities}
The open questions differ in cost and in value. The list below is in a suggested order.
\begin{enumerate}
\item A held-out specimen: fifty unseen words written by hand without the tool. This tests every rule and costs an hour.
\item A decision on the stressed vowel of \emph{words} (\ar{وَرْدْزْ} against the computed \ar{وُرْدْزْ}). For \emph{written} the author chose the computed \ar{رِتِنْ} over the first composed form \ar{رِتَنْ}, which settles the unstressed vowel before a final \emph{n} as kasra.
\item A decision on the open-syllable yeh in words such as \emph{familiar} and \emph{recognise}, where the stressed syllable is closed in speech but open in spelling.
\item A decision on whether to write the stress of English words.
\item A choice of pronunciation source, including regional dictionaries.
\item A context model for the reverse reader.
\item A keyboard layout generated from the tables.
\end{enumerate}

\section{Homographs}
Words such as \emph{read}, \emph{lead} and \emph{live} have more than one pronunciation. The forward transliterator takes the first dictionary entry. A part-of-speech tagger or a word-sense list could choose among them, and the author could instead mark the intended reading by hand. For now, a per-word override handles the cases the author cares about.

\section{Proper names and coinages}
Names the dictionary does not contain are handled by the fallback, whose weak spots appear whenever it logs a word. The override file is the place to record the author's reading of a name. A list of the first hundred such words from the author's own writing would improve the notation on the texts the author actually produces.

\section{Dialects and accents}
The dictionary reflects General American speech. The notation could carry an accent profile, as it already carries profiles for $v$ and hard $g$, and a profile for a British, Canadian or other accent would change the phoneme source and perhaps a few rules. Nothing in the present design prevents this, but no specimen distinguishes them.

\section{Concluding reflection}
The project has a theoretical side that this book has kept apart from its tools. The author connects it to Deacon's account of symbolic reference and to the reading of film in the \emph{Psychocinema} discussion, and neither is needed to use or evaluate the notation. The strongest version of the connection rests on the reader's experience. Reading English in an unfamiliar script makes the normally automatic work of recognition temporarily conspicuous. The reader sees that a word is being identified from a pattern of marks, that the pattern is a convention, and that the convention can be learned and then forgotten again as it becomes fluent. That observation stands on its own and gives the project a philosophical interest. Nothing in the readability of the notation depends on accepting a larger theory of language or of cinema.

\appendix
\chapter{Command Line Reference}\label{app:cli}
The package is run as \texttt{python -m arabglish} with the source directory on the path. The original wrapper \texttt{main.py} is kept for compatibility and takes a text or an input file.

\begin{description}
\item[\texttt{fwd}] English to Arabic script. Options: \texttt{--density} (full, medial, vowels, none), \texttt{--v} (b, f, context), \texttt{--g} (ghain, jeem), \texttt{--persian}, \texttt{--digits} (indic, latin, spoken), \texttt{--attested}, \texttt{--the-before-vowel}, \texttt{--text}, \texttt{-i}, \texttt{-o}.
\item[\texttt{rev}] Arabic script to ranked English candidates.
\item[\texttt{lint}] Check a text against the inventory; \texttt{--fix} rewrites extended letters.
\item[\texttt{stats}] Letter, mark and issue counts for a file.
\item[\texttt{glossary}] Three-line entries (Arabic, respelling, English).
\item[\texttt{roundtrip}] Measure reader recovery over frequent words.
\end{description}

Examples:
\begin{verbatim}
python3 main.py "the nature of lack"
python3 -m arabglish fwd --density none --text "the nature of lack"
python3 -m arabglish lint specimens/atlas.txt
\end{verbatim}
The test suite runs with \texttt{pytest}. Tables for this book are regenerated by \texttt{scripts/make\_tex\_tables.py}.

\section{Sample sessions}
Forward transliteration with different densities:
\begin{verbatim}
$ python3 main.py "The book is written in English."
ذَا بُوكْ إِزْ رِتِنْ إِنْ إِنْغْلِشْ.

$ python3 -m arabglish fwd --density none \
      --text "The book is written in English."
ذا بوك إز رتن إن إنغلش.
\end{verbatim}

Switching the open decisions:
\begin{verbatim}
$ python3 -m arabglish fwd --v b --text "very live"
$ python3 -m arabglish fwd --g jeem --text "good game"
\end{verbatim}

Checking and repairing text:
\begin{verbatim}
$ python3 -m arabglish lint specimens/atlas.txt
$ python3 -m arabglish lint --fix specimens/atlas.txt -o atlas_fixed.txt
\end{verbatim}

Rebuilding the book's tables and passages:
\begin{verbatim}
$ PYTHONPATH=src python3 scripts/make_tex_tables.py
$ PYTHONPATH=src python3 scripts/make_passages.py
$ cd monograph && xelatex main.tex && xelatex main.tex
\end{verbatim}

\chapter{Attested Forms}\label{app:attested}
Forms written by the author in the specimens, with the form computed at the default settings. Source \texttt{P} is the handwritten page, \texttt{T} the typed passage.

{\small\begin{longtable}{@{}lllll@{}}
\toprule
English & Attested & Source & Computed & Match\\
\midrule
\endhead
a & \ar{اَ} & T & \ar{أَ} & differs\\
and & \ar{أَنْدْ} & TP & \ar{أَنْدْ} & exact\\
as & \ar{أَزْ} & T & \ar{أَزْ} & exact\\
both & \ar{بُوثْ} & T & \ar{بُوثْ} & exact\\
built & \ar{بِلْتْ} & T & \ar{بِلْتْ} & exact\\
capitalism & \ar{كَابِتَالِيزْمْ} & T & \ar{كَابِتَلِزْمْ} & differs\\
capitalist & \ar{كَابِتَالِسْتْ} & TP & \ar{كَابِتَلِسْتْ} & differs\\
concept & \ar{كُونْسِبْتْ} & T & \ar{كُونْسِبْتْ} & exact\\
cultural & \ar{كَلْتْشَرَلْ} & TP & \ar{كَلْتْشَرَلْ} & exact\\
delves & \ar{دِلْفْزْ} & TP & \ar{دِلْفْزْ} & exact\\
denial & \ar{دِنَايَلْ} & T & \ar{دِنَايَلْ} & exact\\
explores & \ar{إِكْسْبْلُورْزْ} & T & \ar{إِكْسْبْلُورْزْ} & exact\\
for & \ar{فُورْ} & T & \ar{فُورْ} & exact\\
frames & \ar{فْرَيْمْزْ} & T & \ar{فْرَيْمْزْ} & exact\\
helen & \ar{هِيلِنْ} & TP & \ar{هِيلِنْ} & exact\\
ideology & \ar{آيدِيُولُوجِي} & T & \ar{آيْدِيَالَجِي} & differs\\
in & \ar{إِنْ} & TP & \ar{إِنْ} & exact\\
individual & \ar{إِنْدِبِدْيُوَالْ} & T & \ar{إِنْدَفِجَوَلْ} & differs\\
intersects & \ar{إِنْتِرْسِكْتْسْ} & T & \ar{إِنْتِرْسِكْتْسْ} & exact\\
into & \ar{إِنْتُو} & TP & \ar{إِنْتُو} & exact\\
it & \ar{إِتْ} & T & \ar{إِتْ} & exact\\
its & \ar{إِتْسْ} & TP & \ar{إِتْسْ} & exact\\
lack & \ar{لَاكْ} & TP & \ar{لَاكْ} & exact\\
nature & \ar{نَيْتْشَر} & P & \ar{نَيْتْشَرْ} & letters\\
nature & \ar{نَيْشَرْ} & T & \ar{نَيْتْشَرْ} & differs\\
of & \ar{أُفْ} & TP & \ar{أُفْ} & exact\\
on & \ar{أُونْ} & T & \ar{أُونْ} & exact\\
part & \ar{بَارْتْ} & TP & \ar{بَارْتْ} & exact\\
psychoanalytic & \ar{سَايْكُوأَنَالِيتِيكْ} & T & \ar{سَايْكُوأَنَلِتِكْ} & differs\\
psychocinema & \ar{سَايْكُوسِينِمَا} & TP & \ar{سَايْكُوسِينِمَا} & exact\\
psychology & \ar{سَايْكُولُوجِي} & T & \ar{سَايْكُولُوجِي} & exact\\
ramification & \ar{رَامِفِكَيْشَنْ} & T & \ar{رَامِفِكَيْشَنْ} & exact\\
repression & \ar{رِبْرِشَنْ} & TP & \ar{رِبْرِشَنْ} & exact\\
rollins & \ar{رُولِنْزْ} & TP & \ar{رُولِنْزْ} & exact\\
she & \ar{شِي} & TP & \ar{شِي} & exact\\
social & \ar{سُوشَلْ} & T & \ar{سُوشَلْ} & exact\\
structures & \ar{سْتْرَكْشُورْزْ} & T & \ar{سْتْرَكْتْشَرْزْ} & differs\\
symptoms & \ar{سِمْبْتُمْزْ} & TP & \ar{سِمْبْتُمْزْ} & exact\\
system & \ar{سِسْتَمْ} & T & \ar{سِسْتَمْ} & exact\\
the & \ar{ذَا} & TP & \ar{ذَا} & exact\\
this & \ar{ذِسْ} & TP & \ar{ذِسْ} & exact\\
truth & \ar{تُرُوثْ} & T & \ar{تْرُوثْ} & letters\\
universal & \ar{يُونِبِرْسَلْ} & T & \ar{يُونِفِرْسَلْ} & differs\\
utopianism & \ar{يُتُوبِيَانِيزْمْ} & TP & \ar{يُتُوبِيَنِزْمْ} & differs\\
with & \ar{وِذْ} & T & \ar{وِذْ} & exact\\
\bottomrule
\end{longtable}}

\chapter{Legacy Sources}\label{app:legacy}
Earlier files are preserved unmodified in the project.

\begin{itemize}
\item \texttt{toarabic.ahk}, an AutoHotkey layout from the author's repository. It assigns one Arabic letter to each Latin key, so \texttt{p} and \texttt{v} map to key assignments and not to claims about how the sounds are written.
\item \texttt{pronouncing\_translit\_chat\_version.py}, the pronouncing-based script as pasted from a chat. Its defects are that the vowel table mixes letters and diacritics, the entry for \texttt{sh} is never reached, the voiced-$v$ test indexes letters with a phoneme index, and no vowel or sukoon placement is done.
\item \texttt{cli\_main\_as\_pasted.py}, the command-line wrapper, which imports \texttt{transliterate\_text} from a module named \texttt{transliterator}. A shim of that name forwards to the package.
\end{itemize}

\section{The old version}
The author's earlier program (the files \texttt{old\_version\_}\allowbreak\texttt{transliterator.py} and \texttt{old\_version\_}\allowbreak\texttt{mappings.py}) converts spelling to the first CMU pronunciation and then writes one letter and one mark per phoneme, with nine fixed words. Its consonant choices are $p\to$\ar{ب}, $v\to$\ar{ف}, hard $g\to$\ar{ج}, \emph{ch}$\to$\ar{ش}, \emph{ng}$\to$\ar{ن} and \emph{zh}$\to$\ar{ز}. It merges $p/b$, $v/f$, hard $g$ and $j$, \emph{ch} and \emph{sh}, and \emph{ng} and $n$. It writes each vowel letter before its mark (so AY gives \ar{يَ}), uses no sukoon outside the fixed words, and takes the first dictionary pronunciation without regard to context. Unrecognised phonemes are dropped silently. The author describes it as not definitive.

The current toolkit differs in three ways that matter. It places the vowel mark on the preceding consonant and writes glides with a sukoon (\ar{سَايْ}), which is the form on the handwritten page. It writes hard $g$ as \ar{غ} and \emph{ch} as \ar{تْش} by default, with switches (\texttt{--g jeem}) to return to the old choices. And the $v$ choice is a switch, with \ar{ف} the default as in the old version, since the author regards $v$ as a voiced $f$; \emph{individual} and \emph{universal} in the typed passage use \ar{ب} and are reproduced by \texttt{--v b}. Both versions use only standard Arabic letters, and the author accepts near equivalents such as \ar{ب} for $p$ when exact recovery of short words is not possible.

\end{document}
